\chapter{Ricci flow and curvature estimates}
\label{ch:ricci-flow-v4}

Ricci flow changes a metric according to its Ricci curvature.  Its local
existence theory allows us to start the evolution; curvature estimates
determine how long it continues and what can appear after rescaling near a
singularity.  The distinction between local bounds, bounds on a whole time
slab, and bounds on individual slices is essential throughout this chapter.

All manifolds below have no boundary and are smooth, Hausdorff, and second
countable.  A closed manifold is compact.  We use its Borel measurable
structure and its Riemannian volume.  An interval may include an endpoint;
smoothness and time derivatives there are understood from within the interval.
\begin{definition}[Ricci flow]\label{def:rf-flow}
\lean{PoincareMT.RicciFlow}
\uses{def:foundations-manifold}
A Ricci flow on an interval $J$ is a smooth family of positive definite
metrics $g(t)$ with
\begin{equation}\label{eq:rf-equation}
                         \partial_t g=-2\operatorname{Ric}.
\end{equation}
Its connection is always the Levi--Civita connection of the same metric.
\end{definition}
Completeness means metric completeness on every component, for every
specified time.  On a disconnected manifold the length distance between
different components is infinite.

Our curvature convention is
$\operatorname{Rm}(u,v,w,z)=\langle\mathcal R(u,v)z,w\rangle$,
with $\operatorname{Rm}(u,v,u,v)>0$ on an orthonormal pair on a positively
curved sphere.  We write $R$ for scalar curvature and use the Hilbert tensor
norm for $\operatorname{Rm}$ and its covariant derivatives.  The rough
Laplacian has scalar principal part $g^{ij}\partial_i\partial_j$.
The curvature operator is the self-adjoint operator on $\Lambda^2 TM$
whose quadratic form on $u\wedge v$ is
$\operatorname{Rm}(u,v,u,v)$.  Nonnegative curvature operator means
nonnegativity on every two-form, a stronger condition than nonnegative
sectional curvature in general dimensions.

\section{Scaling and normalized initial data}
\label{sec:rf-normalization}

For a constant $c>0$, the metrics $g$ and $cg$ have the same connection.
Their distance, volume and curvature norms satisfy
\begin{equation}\label{eq:rf-scaling}
 d_{cg}=\sqrt c\,d_g,\qquad d\mu_{cg}=c^{n/2}d\mu_g,\qquad
 |\operatorname{Rm}|_{cg}=c^{-1}|\operatorname{Rm}|_g.
\end{equation}
Indeed the connection formula contains one inverse metric and one
derivative of the metric, so its factors of $c$ cancel.  The covariant
curvature tensor acquires one factor of $c$ and its squared norm four
inverse factors, giving the last identity.  Consequently the parabolic
rescaling
\[
                    \widetilde g(s)=Q\,g(t_*+s/Q),\qquad Q>0,
\]
is again a Ricci flow on precisely those $s$ for which $t_*+s/Q\in J$.

\begin{definition}[Normalized initial data]\label{def:rf-normalized}
\lean{PoincareMT.NormalizedInitialMetric}
\uses{def:foundations-manifold}
A normalized metric on a three-manifold is complete, has finite total
volume, and satisfies, for every $x$ and every $0<r\leq1$,
\[
 |\operatorname{Rm}_{g_0}(x)|\leq1,\qquad
 \mu_{g_0}(B_{g_0}(x,r))\geq\tfrac12\omega_3r^3.
\]
Here $\omega_3$ is the Lebesgue volume of the Euclidean unit ball.
\end{definition}

\begin{proposition}[Existence of normalized metrics]\label{prop:rf-normalize}
\lean{PoincareMT.existsNormalizedInitialMetric}
\uses{def:rf-normalized}
Every closed smooth three-manifold admits a normalized metric.
\end{proposition}
\begin{proof}
Choose a smooth metric $g$.  Compactness gives $B<\infty$ with
$|\operatorname{Rm}_g|\leq B$.  In coordinates, the metric and volume
density converge to their Euclidean values at the center.  A finite
coordinate cover makes the resulting radius uniform: for some $r_0>0$,
$\mu_g(B_g(x,r))\geq\omega_3r^3/2$ for every $x$ and $0<r\leq r_0$.
Choose $c>\max\{B,r_0^{-2}\}$.  Formula~\eqref{eq:rf-scaling} gives the
curvature bound, and for $0<r\leq1$ gives
\[
 \mu_{cg}(B_{cg}(x,r))
 =c^{3/2}\mu_g(B_g(x,r/\sqrt c))\geq\tfrac12\omega_3r^3.
\]
Compactness supplies finite volume and completeness.  When the initial
bound is instead on unit-input curvature evaluations, the elementary
three-dimensional estimate $|\operatorname{Rm}|\leq9B$ permits the
explicit choice $c=\max\{9B,r_0^{-2}\}+1$.
\end{proof}
This is the normalization of \cite[Definition 4.10, p.~67]{MT2007}.
Its compatibility with the logarithmic pinching condition is proved in
Section~\ref{sec:rf-pinching}.

\section{Distance comparison and cut-locus supports}
\label{sec:rf-distance-comparison}

We record the geometric comparison used for escaping limits and
Busemann levels.  All distances and covariant derivatives in this
section belong to one fixed metric.

\begin{lemma}[Squared-distance and triangle comparison]
\label{lem:rf-squared-distance-comparison}
\lean{PoincareMT.RiemannianMetric.squared_distance_sub_sq_concave,PoincareMT.RiemannianMetric.toponogov_hinge,PoincareMT.RiemannianMetric.toponogov_corresponding_side}
\uses{def:foundations-manifold}
Let $(M,g)$ be a complete connected Riemannian manifold with
nonnegative sectional curvature.  For every $p\in M$ and every
unit-speed minimizing geodesic $\eta$, the function
\[
                  t\longmapsto d(p,\eta(t))^2-t^2
\]
is concave on its parameter interval.
If $\alpha:[0,a]\to M$ and $\beta:[0,b]\to M$ are unit-speed
minimizing segments starting at $p$, with $a,b>0$, put
$c=d(\alpha(a),\beta(b))$.  Then
\[
 c^2\leq a^2+b^2-2ab
              \langle\dot\alpha(0),\dot\beta(0)\rangle,
\]
and, for $0\leq s\leq a$ and $0\leq t\leq b$,
\[
 d(\alpha(s),\beta(t))^2
 \geq s^2+t^2-2st\,\frac{a^2+b^2-c^2}{2ab}.
\]
\end{lemma}
\begin{proof}
The index form of a minimizing segment is nonnegative on fields
vanishing at its endpoints. Comparing the endpoint Jacobi field
with a linear parallel field gives
$\operatorname{Hess}(\rho^2)\leq2g$ on a regular inverse-exponential branch.
At a cut point, move the initial point a fraction $\theta$ along a
minimizing segment. The remaining segment has a nonsingular inverse
branch, and the shifted squared radius is an upper support with Hessian
at most $2g/(1-\theta)$. Letting $\theta\downarrow0$ gives approximate
upper supports for the squared distance along $\eta$.

Subtracting $t^2$ leaves supports with arbitrarily small positive
second derivative. A negative minimum below a chord, perturbed by
$\epsilon(t-u)(w-t)$, would have negative second derivative at its
contact point, a contradiction. This proves concavity.
For the hinge inequality, the supporting derivative at the common
initial point is $-2a\langle\dot\alpha(0),\dot\beta(0)\rangle$.
For corresponding sides, apply concavity in each variable to
$d(\alpha(s),\beta(t))^2-s^2-t^2$ and use its values on the two axes.
\end{proof}

The triangle statements are the two forms of
\cite[Theorem 2.4, p. 23]{MT2007}.  The shortened-segment
support and the reversal argument are the Calabi construction
in \cite[Lemma 7.1.9, pp. 284--285]{Petersen2016}.
The proof above keeps the selected minimizing segment and
its inverse branch, as needed when later estimates retain
a specified direction or a fixed convergence map.

\section{Local existence, uniqueness, and continuation}
\label{sec:rf-local}

\begin{theorem}[Local existence, uniqueness, and continuation]\label{thm:rf-local}
\lean{PoincareMT.ricciFlowLocalTheory,PoincareMT.ricciFlowContinuation}
\uses{lem:rf-short-time-family,lem:rf-uniqueness,lem:rf-continuation}
Let $M$ be a closed smooth $n$-manifold and $g_0$ a smooth metric.
There is a $T>0$ and a Ricci flow on $[0,T)$ with initial metric $g_0$.
Two Ricci flows on intervals with least element zero and the same initial
metric agree at every time in the intersection of their intervals.
If $0<T<\infty$ and a flow on $[0,T)$ satisfies
\[
 \exists C\geq0\quad\forall t\in[0,T)\quad\forall x\in M,\qquad
                       |\operatorname{Rm}|(x,t)\leq C,
\]
then it extends, agreeing on $[0,T)$, to $[0,T')$ for some $T'>T$.
\end{theorem}

\subsection{Constructing and comparing flows}

\begin{lemma}[Short-time metric family]\label{lem:rf-short-time-family}
\lean{PoincareMT.exists_ricciFlow_metricFamily}
\uses{def:rf-flow}
For every smooth metric $g_0$ on a closed smooth manifold, there are
$T>0$ and a smooth family of positive metrics $g(t)$ on $[0,T)$ with
$g(0)=g_0$ and $\partial_tg=-2\operatorname{Ric}_{g(t)}$.
\end{lemma}
\begin{proof}
Relative to $g_0$, put
$W^k(\bar g)=\bar g^{ij}(\Gamma(\bar g)^k_{ij}-\Gamma(g_0)^k_{ij})$.
The difference of connections is tensorial. Solve
\begin{equation}\label{eq:rf-deturck}
 \partial_t\bar g=-2\operatorname{Ric}_{\bar g}
                 +\mathcal L_{W(\bar g)}\bar g,\qquad \bar g(0)=g_0 .
\end{equation}
Its principal part is $\bar g^{ab}\partial_a\partial_b\bar g_{ij}$.
The Lean construction first solves the linearized equation in the
spectral $L^2$ scale of a positive elliptic operator. Tensor probes,
integration by parts, and compact inclusion give the closed generator
and its resolvent. Splitting the nonlinear residual into a small
highest-derivative coefficient and lower-order products gives a
contraction on a ball of radius $\rho$ for a sufficiently short time.
The order of choices is $\rho$, then the lifetime; the resulting metric
stays positive, for example $\bar g\geq 3g_0/4$.

Leibniz estimates use $L^\infty$ control of lower jets and $L^2$ control
of the highest jets. Smooth dependence on localized diffeomorphism
parameters and covariance identify the parameter orbit with the actual
pullback orbit, yielding spatial smoothness; the equation then gives
the time derivatives. Finally solve
$\partial_t\phi=-W(\bar g)\circ\phi$ and set $g=\phi^*\bar g$.
Differentiating the pullback cancels the Lie derivative in
\eqref{eq:rf-deturck}. The negative sign in the gauge ODE is essential.
Choosing the Levi--Civita connection of each metric assembles the
family into Definition~\ref{def:rf-flow}.
\end{proof}

\begin{lemma}[Uniqueness from difference energy]\label{lem:rf-uniqueness}
\lean{PoincareMT.ricciFlowUniqueness}
\uses{def:rf-flow}
Two Ricci flows on a closed smooth manifold, whose intervals have
least element zero and whose metrics agree at zero, agree on the
intersection of their time intervals.
\end{lemma}
\begin{proof}
Compare the metrics, connections, and curvature tensors simultaneously.
On a compact subinterval their coefficients and derivatives are bounded.
The squared differences define a nonnegative integrated energy $E$;
integration by parts in the curvature equation produces a dissipative
gradient term $D$. The three energy estimates have the form
$CE$, $CE+D/2$, and $CE-D$. Summing absorbs the gradient terms and
gives $E'\leq CE$. Since $E(0)=0$, Gronwall gives $E=0$.
Exhaust the common time interval by compact subintervals.
\end{proof}

\begin{lemma}[Continuation through a bounded-curvature endpoint]
\label{lem:rf-continuation}
\lean{PoincareMT.ricciFlowContinuation}
\uses{lem:rf-short-time-family,lem:rf-uniqueness}
If a Ricci flow on a closed smooth manifold is defined on $[0,T)$,
where $0<T<\infty$, and one constant bounds $|\operatorname{Rm}|$
everywhere on this interval, it extends smoothly beyond $T$.
\end{lemma}
\begin{proof}
The local estimates in the continuation construction bound all spatial
metric jets away from zero. The evolution equation bounds their time
derivatives, so the jets converge as $t\uparrow T$. Metric comparison
keeps the limiting metric positive. These compatible limits produce a
smooth terminal metric $g_T$.
Apply Lemma~\ref{lem:rf-short-time-family} at $g_T$ and translate its
initial time to $T$. The equation identifies the time jets recursively
with spatial jets, so the old and new solutions glue smoothly.
Lemma~\ref{lem:rf-uniqueness} gives agreement on overlaps.
\end{proof}

These three constructions prove Theorem~\ref{thm:rf-local}; compare
\cite[Theorem 3.11, pp.~39--40; Proposition 4.12, p.~68]{MT2007}.
The implementation of the local theory constructs its own analytic
estimates; it does not assume the later assembled curvature theory.

\section{Evolution, derivative estimates, and scalar comparison}
\label{sec:rf-curvature}

For fixed tangent inputs, put
$B(u,v,w,z)=\sum_{a,b}\operatorname{Rm}(u,e_a,v,e_b)
                         \operatorname{Rm}(w,e_a,z,e_b)$
in an orthonormal basis.  Define
\[
\begin{split}
 Q(u,v,w,z)={}&2\{B(u,v,w,z)-B(u,v,z,w)\\
              &\hspace{13mm}-B(u,z,v,w)+B(u,w,v,z)\}\\
 &-\operatorname{Rm}(\operatorname{Ric}^{\sharp}u,v,w,z)
  -\operatorname{Rm}(u,\operatorname{Ric}^{\sharp}v,w,z)\\
 &-\operatorname{Rm}(u,v,\operatorname{Ric}^{\sharp}w,z)
  -\operatorname{Rm}(u,v,w,\operatorname{Ric}^{\sharp}z).
\end{split}
\]
These are contractions of tensors, hence basis independent.

\begin{proposition}[Evolution identities]\label{prop:rf-evolution}
\lean{PoincareMT.RicciFlow.hasDerivWithinAt_scalarCurvature,PoincareMT.RicciFlow.hasDerivWithinAt_curvatureTensor,PoincareMT.RicciFlow.hasDerivWithinAt_ricci}
\uses{def:rf-flow}
On every smooth Ricci flow, scalar curvature is smooth in space and time,
and
\[
\begin{split}
 (\partial_t-\Delta)\operatorname{Rm}&=Q,\\
 (\partial_t-\Delta)\operatorname{Ric}(u,v)
 &=2\sum_{a,b}\operatorname{Rm}(u,e_a,v,e_b)\operatorname{Ric}(e_a,e_b)
       -2\sum_a\operatorname{Ric}(u,e_a)\operatorname{Ric}(e_a,v),\\
 (\partial_t-\Delta)R&=2|\operatorname{Ric}|^2.
\end{split}
\]
At an included endpoint these are derivatives within the time interval.
Also $|\nabla^0\operatorname{Rm}|=|\operatorname{Rm}|$.
\end{proposition}
\begin{proof}
Differentiating metric compatibility and torsion freeness gives
\[
 \partial_t\Gamma^k_{ij}
 =-g^{k\ell}(\nabla_i\operatorname{Ric}_{j\ell}
             +\nabla_j\operatorname{Ric}_{i\ell}
             -\nabla_\ell\operatorname{Ric}_{ij}).
\]
Insert this in the antisymmetrized derivative defining curvature.
Commuting the two covariant derivatives contributes the curvature
action on each tensor slot.  The differential Bianchi identity turns
the remaining second derivatives into the rough Laplacian; the
quadratic contractions group into the four displayed $B$ terms.
Differentiating the lowering metric supplies the fixed-input corrections
in $Q$.  In each trace one must also differentiate the inverse metric:
$\partial_tg^{ij}=2\operatorname{Ric}^{ij}$.  The first trace gives the
Ricci formula and the second cancels its negative Ricci-square term,
leaving $2|\operatorname{Ric}|^2$.  These calculations and the curvature
symmetries also prove the claimed regularity.
\end{proof}
The fixed-input convention matters in these contractions; compare
\cite[equations (3.5)--(3.6), p.~41; Lemma 3.15, p.~42]{MT2007}.

\begin{theorem}[Local derivative estimates]\label{thm:rf-shi}
\lean{PoincareMT.local_curvatureDerivative_bound,PoincareMT.local_curvatureDerivative_bound_of_initial}
\uses{prop:rf-evolution}
For every $n,k\in\mathbb N$ and $K,\alpha,r>0$ there is a constant
$C=C(n,k,K,\alpha,r)>0$ with the following property.  Let $0<T\leq\alpha/K$
and let $g(t)$ be a Ricci flow on $[0,T]$.  If
$\overline{B_{g(0)}(p,r)}$ is compact and $|\operatorname{Rm}|\leq K$ on
$B_{g(0)}(p,r)\times[0,T]$, then
\[
 |\nabla^k\operatorname{Rm}|(x,t)\leq C t^{-k/2}
 \quad (x\in B_{g(0)}(p,r/2),\ 0<t\leq T).
\]
There is also, for every $n,k,l,K,\alpha,r$, a constant with the same
independence from the manifold, center and flow such that
\[
 |\nabla^k\operatorname{Rm}|(x,t)
       \leq C t^{-\max\{k-l,0\}/2}.
\]
For this second assertion assume the same compact-ball condition,
$|\operatorname{Rm}|\leq K$ on the whole manifold throughout $[0,T]$,
and $|\nabla^j\operatorname{Rm}|(x,0)\leq K$ on the whole manifold for
$0\leq j\leq l$.  The conclusion holds on the half ball for
$t\in[0,T]$ with $t>0$ or $k\leq l$; when $k\leq l$ its right side is $C$,
including at zero.
\end{theorem}

\begin{proof}
Write $N_j=|\nabla^j\operatorname{Rm}|$. Differentiating the
curvature equation and commuting covariant derivatives gives
\[
 (\partial_t-\Delta)N_k^2
 \leq-2N_{k+1}^2+
 C(n,k)\sum_{i+j=k}N_iN_jN_k.
\]
Localize on nested balls inside the fixed initial ball. The
curvature bound compares the time-dependent metric with the initial
one, so the cutoff remains supported in a compact spatial cylinder.
At cut points, shortened minimizing segments give smooth upper
supports for distance; the maximum principle applies to these
supports and then passes to the limiting cutoff.

Induct on $k$. A Bernstein expression coupling $N_k^2$ with the
already bounded lower derivatives retains a negative multiple of
$N_k^4$, which absorbs the cutoff-gradient errors at a positive
maximum. Parabolic rescaling of half-time windows gives
$t^{k/2}N_k\leq C$. The constants are chosen from
$n,k,K,\alpha,r$ before the manifold, flow, and center.

When derivatives through order $l$ are initially bounded everywhere,
the same induction starts from those bounds and only pays a factor
$t^{-1/2}$ for each derivative above $l$. This gives
$t^{-\max(k-l,0)/2}$. The initial-time case is retained only for
$k\leq l$; otherwise the asserted bound is at positive times.
\end{proof}

Compare \cite[Theorems 3.28--3.29 and their proof, pp.~51--57]{MT2007}.
The product-energy argument follows that proof.  Here the scale is
fixed on each half-time interval, and the cutoff is built from evolving
distance supports on compact carriers.  This supplies the localization
without requiring the exponential map to be injective on the original
ball.

\begin{proposition}[Metric comparison]\label{prop:rf-metric-comparison}
\lean{PoincareMT.RicciFlow.metric_comparison_of_curvature_bound}
\uses{def:rf-flow}
If $s\leq t$ belong to the flow interval and
$|\operatorname{Rm}|(x,\tau)\leq K$ for all $x$ and $\tau\in[s,t]$, where
$K\geq0$, then
\[
 e^{-2nK(t-s)}g(s)\leq g(t)\leq e^{2nK(t-s)}g(s).
\]
\end{proposition}
\begin{proof}
Contraction gives $|\operatorname{Ric}(v,v)|\leq nK g(v,v)$.
For a fixed nonzero tangent vector,
$|(\log g(\tau)(v,v))'|\leq2nK$.  Integrate from $s$ to $t$;
the zero vector case follows directly.  The same argument is local in
space if a curvature bound is known only along the fixed point.
\end{proof}

\begin{proposition}[Scalar lower barrier]\label{prop:rf-scalar}
\lean{PoincareMT.RicciFlow.scalarCurvature_lowerBound,PoincareMT.RicciFlow.scalarCurvature_lowerBound_three}
\uses{prop:rf-evolution}
Let $n>0$ and let $g(t)$ be a Ricci flow on a closed $n$-manifold on
$[a,b)$, where $a<b$.  If $r_0<0$ and $R(x,a)\geq r_0$ for all $x$, then
\[
 R(x,t)\geq \frac{r_0}{1-2r_0(t-a)/n}\quad(a\leq t<b).
\]
In dimension three, if $a\geq0$ and $R(x,a)\geq-6/(1+4a)$, then
$R(x,t)\geq-6/(1+4t)$ on $[a,b)$.
\end{proposition}
\begin{proof}
Cauchy--Schwarz for the Ricci eigenvalues gives
$|\operatorname{Ric}|^2\geq R^2/n$.  The displayed comparison function
$u$ solves $u'=2u^2/n$ and $u(a)=r_0$; its denominator stays positive.
On each compact subinterval,
\[
 (\partial_t-\Delta)(R-u)\geq \tfrac2n(R+u)(R-u).
\]
The coefficient is bounded there.  An exponential weight and a small
strict barrier rule out a first negative minimum, since the Laplacian
is nonnegative at that minimum.  Let the strict perturbation tend to
zero.  Every $t<b$ lies in such a subinterval.  Substitution of
$n=3$ and $r_0=-6/(1+4a)$ gives the last formula.
\end{proof}

\begin{proposition}[Positivity and rigidity]\label{prop:rf-positivity}
\lean{PoincareMT.RicciFlow.nonnegativeSectionalCurvature_preserved,PoincareMT.RicciFlow.nonnegativeRicciCurvature_preserved,PoincareMT.RicciFlow.flat_of_scalarCurvature_eq_zero}
\uses{prop:rf-evolution}
On a closed three-manifold, a Ricci flow on $[0,T]$, $T>0$, preserves
nonnegative sectional curvature and preserves nonnegative Ricci curvature.
On a connected three-manifold, a flow on $[0,T]$ with nonnegative
sectional curvature at every time and $R(p,T)=0$ has
$|\operatorname{Rm}|(x,t)=0$ for every $x$ and $0\leq t\leq T$.
The rigidity assertion requires no compactness or completeness hypothesis.
\end{proposition}
\begin{proof}
In an evolving orthonormal frame the curvature operator eigenvalues
$\lambda,\mu,\nu$ in dimension three have reaction
$2(\lambda^2+\mu\nu)$ and its cyclic permutations.
At a zero least eigenvalue with the other two nonnegative, its reaction
is nonnegative.  The compact tensor maximum principle therefore
preserves the cone of nonnegative curvature operators, which in
dimension three is the sectional-curvature cone.

For Ricci, suppose $\mu+\nu=0$ is the smallest Ricci eigenvalue.
Then $\lambda+\mu,\lambda+\nu\geq0$, and the reaction of $\mu+\nu$ is
$2(\mu^2+\nu^2+\lambda(\mu+\nu))=4\mu^2\geq0$.
The same cone argument applies.  For rigidity, the scalar evolution
makes $R\geq0$ a supersolution of the heat equation.  The local strong
maximum principle, propagated through connected coordinate
neighborhoods, makes it zero at every earlier positive time if it
vanishes at $(p,T)$.  Continuity includes both endpoints.  A
nonnegative curvature operator in dimension three has nonnegative
eigenvalues whose sum is $R/2$; they therefore all vanish.
\end{proof}
See \cite[Proposition 4.1, pp.~64--65; Corollaries 4.14--4.15,
pp.~69--70; Theorem 4.18, pp.~71--72]{MT2007}.

\section{The Hamilton--Ivey invariant region}
\label{sec:rf-pinching}

In dimension three put
\[
 k_{\min}(x,t)=\inf\{\operatorname{Rm}(u,v,u,v):
                   u,v\text{ orthonormal at }(x,t)\},\qquad
 X(x,t)=\max\{-k_{\min}(x,t),0\}.
\]
The three eigenvalues $k_1\geq k_2\geq k_3$ of the curvature operator obey
\begin{equation}\label{eq:rf-spectrum}
 k_{\min}=k_3,\qquad R=2(k_1+k_2+k_3),\qquad
 |\operatorname{Rm}|^2=4(k_1^2+k_2^2+k_3^2).
\end{equation}
One obtains this operator without choosing an orientation by using
$\tfrac12R\,g-\operatorname{Ric}$ on the tangent bundle; in an oriented
orthonormal frame the Hodge identification with $\Lambda^2$ gives the
same spectrum.  Changing orientation does not change these identities.

\begin{definition}[Hamilton--Ivey pinching]\label{def:rf-pinched}
\lean{PoincareMT.HamiltonIveyPinchedAt}
\uses{def:rf-flow}
A slice is Hamilton--Ivey pinched at absolute time $t$ if $t\geq0$,
$R\geq-6/(1+4t)$, and, at every point with $X>0$,
\begin{equation}\label{eq:rf-log-pinching}
                  R\geq2X\{\log X+\log(1+t)-3\}.
\end{equation}
\end{definition}

\begin{theorem}[Preservation of Hamilton--Ivey pinching]\label{thm:rf-pinching}
\lean{PoincareMT.hamiltonIveyPinching,Poincare.fullNorm_le_of_hamiltonIvey}
\uses{def:rf-pinched,prop:rf-evolution,prop:rf-scalar}
Let $0\leq a<b$ and let $g(t)$ be a Ricci flow on a closed
three-manifold on $[a,b)$.  If its initial slice is pinched at time $a$,
then every slice is pinched at its own time $t$.
At each point where $R\leq R_0$,
\begin{equation}\label{eq:rf-full-bound}
             |\operatorname{Rm}|\leq13\max\{R_0,e^4\}.
\end{equation}
In particular normalized initial data satisfy the initial pinching
condition at time zero.
\end{theorem}

\subsection{Convexity and the reaction calculation}

Write $S=R/2$ and $f_t(X)=X(\log X+\log(1+t)-3)$, continuously extended
by $f_t(0)=0$.  There is a unique $b_*>e^2$ with
$b_*(\log b_*-3)=-3$: strict convexity, the values at $1,e^2,e^3$,
and monotonicity beyond $e^2$ prove this.  Set $b(t)=b_* /(1+t)$.
The convex region used in the tensor argument is
\[
 \mathcal C_t=\{B=B^*:\
       \operatorname{tr}B\geq f_t(\max\{b(t),X(B)\})\},\qquad
 X(B)=\max\{-\lambda_{\min}(B),0\}.
\]
The clipped function is constant $-3/(1+t)$ below $b(t)$ and convex
increasing above it.  Since $\lambda_{\min}$ is the infimum of the
linear Rayleigh functionals, $X(B)$ is convex.  Composition with the
increasing convex clipped function proves convexity of $\mathcal C_t$.
It is closed and orthogonally invariant.

For an ordered spectrum, membership implies $S\geq f_t(X)$ even in the
clipped range.  If $0<X\leq1/(1+t)$, then $f_t(X)\leq-3X\leq S$.
Between $1/(1+t)$ and $b(t)$, convexity puts $f_t(X)$ below the
horizontal level $-3/(1+t)$.  Beyond $b(t)$ the inequality is part of
the definition.  Conversely the stronger scalar bound
$S\geq-3/(1+4t)$ together with \eqref{eq:rf-log-pinching} implies
membership for $t\geq0$.

The key algebra can be checked without differentiating eigenvectors.
On a reaction trajectory diagonalized initially, the fixed eigenbasis
remains diagonal.  Write its ordered entries as
$\lambda\geq\mu\geq\nu$, put $X=-\nu>0$, $Y=-\mu$, and first use the
reaction with rate one.  Then
\[
 S'=X^2+Y^2+\lambda^2+XY-\lambda(X+Y),\qquad
 X'=-X^2+\lambda Y.
\]
Direct expansion gives
\[
 XS'-(S+X)X'=X^3+
             XY^2+\lambda Y(Y-X)+\lambda^2(X-Y).
\]
The last three terms are nonnegative.  If $Y\leq0$, each term can be
written as a product of nonnegative factors, since $\lambda\geq-Y$.
If $Y\geq0$, use the identity
\[
 XY^2+\lambda Y(Y-X)+\lambda^2(X-Y)
 =Y^3+(X-Y)\{(\lambda-Y/2)^2+3Y^2/4\}.
\]
It follows that $(S/X-\log X)'\geq X$ for rate one, and at least $2X$
for the geometric reaction of rate two.  On the increasing logarithmic
boundary $X\geq b(t)>1/(1+t)$, this exceeds the derivative of
$\log(1+t)-3$.  On the horizontal boundary,
$S'=2(\lambda^2+\mu^2+\nu^2+\lambda\mu+\mu\nu+\nu\lambda)
\geq4S^2/3$, so the vector field also points inward.
Repeated eigenvalues are handled by the symmetric polynomial equation
or supporting functionals, without requiring a smooth ordered
eigenbasis.  This proves reaction invariance.

\subsection{From the reaction to the geometric equation}

Transport $\tfrac12R\,g-\operatorname{Ric}$ to a fixed initial Euclidean
bundle using the time-dependent isometries induced by the metric
equation.  Explicitly, at each point solve
$\partial_tP_t=\operatorname{Ric}_{g(t)}^{\sharp}P_t$, $P_a=I$.
Differentiating $g(t)(P_tv,P_tw)$ gives zero: the two input
derivatives cancel $\partial_tg=-2\operatorname{Ric}$.  Thus $P_t$
is an isometry from the initial fiber to the time-$t$ fiber.
Denote the pulled-back endomorphism by $B$.  The moving-frame
evolution has reaction $2(B^2+\operatorname{adj}B)$, with the transported
rough Laplacian as spatial part.  Scaling $U=(1+t)B$ replaces
$\mathcal C_t$ by the fixed set $\mathcal C_0$ and gives
\[
 (\partial_t-\Delta)U
       =(1+t)^{-1}\{U+2(U^2+\operatorname{adj}U)\}.
\]
For every outward supporting functional $\ell$ at $A\in\mathcal C_0$,
reaction invariance implies that $\ell$ of this right side is
nonpositive at $A$.  Indeed start the finite-dimensional reaction
curve at $A$, rescale it to stay in $\mathcal C_0$, apply $\ell$,
and differentiate the resulting one-sided maximum at zero.

Here $\Delta U$ denotes the pullback of the geometric rough
Laplacian, with the same spatially constant factor $1+t$.
Its contact sign requires a spatial argument.  Fix a time and a
point $p$ at which the distance of a smooth tensor $T(x)$ to the
fiberwise carrier $K_x$ has a local maximum.  Choose an active unit
support $(A,N)$ at $p$, meaning
\[
 A\in K_p,\qquad |N|=1,\qquad
 \langle N,Z-A\rangle\leq0\quad(Z\in K_p),\qquad
 \langle N,T(p)-A\rangle=\operatorname{dist}(T(p),K_p).
\]
Extend an orthonormal basis at $p$ by parallel transport along
coordinate rays in a small centered chart.  In coordinates each
extended vector $Y$ solves
\[
 \frac{d}{ds}Y(sz)=-\Gamma(sz)(z,Y(sz)),\qquad Y(0)=v.
\]
Smooth dependence for this linear ODE gives smooth local fields.
Uniqueness identifies the endpoint field on every ray with the
parallel solution along that ray.  At the center this proves
$\nabla_wY=0$ for every $w$; differentiating the same ray equation
once more gives $\nabla_w\nabla_wY=0$ there, where the direction
is extended with constant coordinate components.  The correction
$\nabla_{\nabla_ww}Y$ also vanishes because all first derivatives
vanish.  Only these diagonal second derivatives enter the rough
Laplacian.

Metric compatibility makes the resulting maps $Q_x:T_pM\to T_xM$
isometries.  They preserve $K$, because its definition is invariant
under orthogonal changes of frame.  Write $Q_x$ also for the induced
tensor isometry and set $f(x)=\langle Q_xN,T(x)\rangle$.
The support inequality and Cauchy--Schwarz give
\[
 f(x)-\langle N,A\rangle
   =\langle Q_xN,T(x)-Q_xA\rangle
   \leq\operatorname{dist}(T(x),K_x).
\]
There is equality at $p$, so $f$ has a local maximum at $p$.
Expand $f$ as a finite sum of evaluations of $T$ on the transported
basis fields.  The first and diagonal second derivative identities
remove every input-field term in its Laplacian.  Consequently
\[
 \langle N,\Delta T(p)\rangle=\Delta f(p)\leq0.
\]
Apply this to $(1+t)(\tfrac12Rg-\operatorname{Ric})$ on the
time-$t$ slice.  Isometric pullback by $P_t$ preserves distances,
active supports and this pairing, yielding the required sign for
the diffusion term of $U$.  This transfer uses the slice calculation;
it does not differentiate $P_t$ in space.

To finish comparison on a closed time slab, choose finitely many
compact coordinate neighborhoods with local isometries preserving
the carrier.  Continuity bounds $|U|$ by some $H$ on their union.
Transporting one carrier point from each coordinate fiber also
gives carrier points of uniformly bounded norm, say at most $J$.
The nearest carrier point to $U$ then has norm at most $2H+J$:
its distance from $U$ is at most $H+J$.  Therefore all positive
distances are realized by supports $(A,N)$ with
$|A|\leq2H+J$ and $|N|=1$.  In each coordinate fiber these supports
form a closed bounded set.  The finite union over the compact
neighborhoods is a compact parameter space for the functions
\[
 F_{x,A,N}(t)=\langle N,U(x,t)-A\rangle.
\]
A positive maximum among these functions equals the maximum
distance to the carrier.  At such a maximum the diffusion pairing
is nonpositive by the preceding argument.  If $\psi_t$ denotes
the scaled reaction and $D$ is its uniform Lipschitz constant on
a ball containing both $U$ and its nearest carrier point, inwardness
gives the following estimate.  Put $A_*=\operatorname{proj}_{K_x}U$.
Activity and the support inequality imply
$\langle N,A_*-A\rangle=0$: indeed
$\langle N,U-A_*\rangle\leq|U-A_*|$
and $\langle N,A_*-A\rangle\leq0$, while their sum is exactly
$|U-A_*|$.  Thus $N$ also supports the carrier at $A_*$, and
\[
 \frac{d}{dt}F_{x,A,N}(t)
 \leq\langle N,\psi_t(U)-\psi_t(A_*)\rangle
 \leq D F_{x,A,N}(t).
\]
The last inequality uses $|U-A_*|=F_{x,A,N}(t)$.
Were any $F$ positive, the continuous compact family
$e^{-(D+1)(t-a)}F(t)$ would attain a positive spacetime maximum
at a time greater than $a$.  Its time derivative would be
nonnegative there, including the left derivative at the slab's
right endpoint, whereas the displayed inequality makes it strictly
negative.  This contradiction proves $U\in\mathcal C_0$ without
differentiating the maximum-distance function.
This is the tensor maximum principle of
\cite[Theorem 4.8, p.~66]{MT2007}, applied to the specific carrier just
constructed; the ODE alone would not establish the spatial assertion.
The scalar barrier in Proposition~\ref{prop:rf-scalar} supplies the
stronger scalar component of Definition~\ref{def:rf-pinched}.
Apply the argument on each $[a,t]$ with $t<b$.  The algebra and carrier
are those of \cite[Claims 4.28--4.31, pp.~77--79; Theorem 4.32,
pp.~79--80]{MT2007}.

To prove \eqref{eq:rf-full-bound}, let $K=\max\{R_0,e^4\}$.
If $X>e^4$, the logarithmic inequality and $t\geq0$ imply
$R\geq2X$, hence $X\leq K$; if $X\leq e^4$ this is immediate.
Thus $k_i\geq-K$, while $k_1+k_2+k_3=R/2\leq K/2$ gives
$k_i\leq5K/2\leq3K$.  Formula~\eqref{eq:rf-spectrum} gives
$|\operatorname{Rm}|^2\leq108K^2\leq169K^2$, proving the stated
continuous bound.  Finally, $|\operatorname{Rm}|\leq1$ implies
$X\leq1$ and $R\geq-6$.  For $X>0$, $R/2\geq-3X$ and
$\log X\leq0$, giving the time-zero logarithmic inequality.

\section{Differential Harnack inequalities}
\label{sec:rf-harnack}

\begin{theorem}[Differential and integrated Harnack inequalities]\label{thm:rf-harnack}
\lean{PoincareMT.differentialHarnackAncientTheory}
\uses{prop:rf-bounded-harnack,lem:rf-slab-bound}
Let $g(t)$ be a Ricci flow on an $n$-manifold for $T_0<t<T_1$.
Assume every slice is complete, its curvature operator is nonnegative,
and it is bounded on that slice: for every $t$ there is $K_t\geq0$
bounding its operator norm at every point.  Then for every $(x,t)$
and $v\in T_xM$,
\begin{equation}\label{eq:rf-harnack}
 \partial_tR+\frac{R}{t-T_0}+2\,dR(v)+2\operatorname{Ric}(v,v)\geq0.
\end{equation}
If $T_0<t_1<t_2<T_1$ and $\gamma:[t_1,t_2]\to M$ is $C^1$, put
$E(\gamma)=\int_{t_1}^{t_2}|\gamma'(s)|_{g(s)}^2\,ds$.  Then
\begin{equation}\label{eq:rf-harnack-path}
 R(\gamma(t_2),t_2)(t_2-T_0)\geq
 R(\gamma(t_1),t_1)(t_1-T_0)e^{-E(\gamma)/2}.
\end{equation}
\end{theorem}

\subsection{The Hamilton block and its scalar trace}

For elapsed time $\tau>0$, define tensors
\[
\begin{split}
 P_{ijk}&=\nabla_i\operatorname{Ric}_{jk}
                         -\nabla_j\operatorname{Ric}_{ik},\\
 M_{ij}&=\Delta\operatorname{Ric}_{ij}
       -\tfrac12\nabla_i\nabla_jR
       +2\sum_{a,b}\operatorname{Rm}_{aibj}\operatorname{Ric}_{ab}
       -\sum_a\operatorname{Ric}_{ia}\operatorname{Ric}_{aj}
       +\frac{\operatorname{Ric}_{ij}}{2\tau}.
\end{split}
\]
They form a quadratic form on $\Lambda^2 TM\oplus TM$ with curvature
block $\operatorname{Rm}$, off-diagonal block $P$, and vector block
$M$.  With the two-form convention fixed by decomposable forms, its
value at $(v\wedge w,w)$ is
$\operatorname{Rm}(v,w,v,w)+2P(v,w,w)+M(w,w)$.
Hamilton block positivity says this full quadratic form is nonnegative.

The contracted Bianchi identity and the scalar evolution give
\[
 \sum_iP(v,e_i,e_i)=\tfrac12dR(v),\qquad
 \operatorname{tr}M=\tfrac12\Delta R+|\operatorname{Ric}|^2
                                      +R/(2\tau).
\]
Sum twice the block inequality at $(v\wedge e_i,e_i)$ over $i$.
Since $\sum_i\operatorname{Rm}(v,e_i,v,e_i)=\operatorname{Ric}(v,v)$,
the result is exactly \eqref{eq:rf-harnack}.

\subsection{Reaction, contact, and the noncompact barrier}

\begin{proposition}[Hamilton positivity on a bounded slab]\label{prop:rf-bounded-harnack}
\lean{Poincare.RicciFlow.Harnack.hamiltonBlockPos_on_bounded_slab,Poincare.RicciFlow.Harnack.ancient_differential_of_bounded_curvature}
\uses{prop:rf-evolution,thm:rf-shi}
Let $g(t)$, $a<t<b$, be a Ricci flow with complete slices and
nonnegative curvature operator.  Suppose there is one $K\geq0$ such
that $|\operatorname{Rm}|(x,t)\leq K$ for all $x$ and $a<t<b$.
Then its Hamilton block with elapsed time $t-a$ is positive semidefinite.
In particular \eqref{eq:rf-harnack} holds with $T_0=a$.
For a complete ancient flow with nonnegative curvature operator and
one full-curvature bound on its whole time range, sending $a$ to
$-\infty$ gives
$\partial_tR+2dR(v)+2\operatorname{Ric}(v,v)\geq0$.
Neither conclusion uses an asymptotic-volume-ratio theorem.
\end{proposition}

\begin{proof}
Buffer the initial time and use Theorem~\ref{thm:rf-shi} to bound the
first two curvature derivatives. The Hamilton block then satisfies
a tensor evolution equation whose reaction preserves its positive
cone. At a null vector, the reaction is a sum of squares. This
calculation uses the connection on $\Lambda^2TM\oplus TM$, including
the mixed terms from $P$; entrywise positivity of the diagonal blocks
would not suffice.

For a compact domain, add a strictly positive tensor perturbation.
At a first null contact the scalarized block has nonnegative spatial
Hessian, while its perturbed evolution is strictly positive, excluding
that contact. On a complete noncompact component, construct a proper
smooth exhaustion from the buffered bounded geometry. Its gradient,
Hessian, and time-variation estimates control the error terms in a
growing spatial perturbation. First work on compact sublevel sets;
then send the boundary to infinity and the perturbation to zero.
The reciprocal elapsed-time term in $M$ supplies the initial barrier.

This proves Hamilton positivity at every interior point of the
buffered slab. Move the buffered initial time back to $a$ to recover
the correction $\operatorname{Ric}/(2(t-a))$. Taking the scalar trace
gives the differential Harnack inequality. For a uniformly bounded
ancient flow, send $a\to-\infty$; the elapsed-time correction vanishes.
Neither this limit nor the bounded-slab construction uses the
asymptotic volume ratio.
\end{proof}

\subsection{Why bounds on separate slices suffice}

The finite Harnack theorem assumes a separate bound on each slice.
Pointwise continuity in time does not make those bounds uniform on
a time slab. The proof first obtains a uniform initial ball volume,
then propagates it, and only then upgrades curvature control.

\begin{lemma}[Initial component volume]\label{lem:rf-initial-volume}
\lean{PoincareMT.RicciFlow.exists_initial_unit_ball_volume_lower_bound_on_component}
\uses{lem:rf-squared-distance-comparison,prop:rf-evolution}
On a complete slice of dimension $n\geq1$, suppose the curvature
operator is nonnegative and bounded on a connected component $N$.
There is $v>0$ such that $\mu(B(x,1))\geq v$ for all $x\in N$.
The constant may depend on the entire component.
\end{lemma}
\begin{proof}
The bounded sectional curvature gives a uniform nonsingularity
radius for exponential maps by Jacobi-field comparison.
Construct a proper convex Busemann exhaustion $b$ and consider
the injectivity radius truncated below that nonsingularity radius.
It is positive and lower semicontinuous, so has a positive minimum
on the compact zero level of $b$.

If a smaller injectivity radius occurred elsewhere, minimize
$i_c+\eta b$ on a compact sublevel. The first colliding exponential
branches have opposite terminal velocities. Moving their common
endpoint toward $\{b=0\}$ therefore changes their average length
only to second order, while decreasing $\eta b$ to first order.
This contradicts minimality. Hence the injectivity lower bound is
global. Jacobi comparison controls the Jacobian on a uniform
exponential ball; change of variables gives the claimed volume bound.
\end{proof}
Compare \cite[Lemmas 3.8 and 3.13, pp.~15--19;
Lemma 5.1, pp.~35--37]{MaederBaumdickerSeidel2025}.

\subsection{Scalar integrals and regular fibers}
\label{sec:rf-scalar-integral}

\begin{theorem}[Uniform unit-ball scalar integral]\label{thm:rf-scalar-integral}
\lean{PoincareMT.exists_uniform_unitBall_scalar_integral_bound}
\uses{lem:rf-squared-distance-comparison}
For each dimension $n$ there is $P_n>0$ such that every complete
$n$-manifold $(N,h)$ with $\sec_h\geq-1$ satisfies
\[
             \int_{B_h(p,1)}R_h\,d\mu_h\leq P_n
                         \quad\hbox{for every }p\in N.
\]
It is Petrunin's estimate \cite[Theorem 1.1, manuscript p.~1]{Petrunin2009}.
\end{theorem}
\begin{proof}
On a compact regular slab $f^{-1}([u,v])$, write
$\mathbf n=\nabla f/|\nabla f|$, let $B$ be the second fundamental
form of the levels, and set $H=\operatorname{tr}B$ and $G=H^2-|B|^2$.
Integration by parts, followed by coarea, gives
\[
 \int_{f^{-1}([u,v])}(G-\operatorname{Ric}(\mathbf n,\mathbf n))\,d\mu
 =\int_{f^{-1}(v)}H\,dA-\int_{f^{-1}(u)}H\,dA.
\]
Both endpoint normals point toward increasing $f$. Together with
$R_{f^{-1}(s)}=R-2\operatorname{Ric}(\mathbf n,\mathbf n)+G$,
this reduces an ambient scalar integral to level integrals and
controlled boundary terms.

The induction is on fiber dimension, simultaneously for every
codimension. Smooth approximations to distance on compact annuli
supply regular fibers with quantitative gradient and Hessian bounds.
An additional almost-opposite pair of gradients raises the
codimension by one. The induced Gauss formula retains the
second-fundamental-form error. Coarea, area variation, and volume
comparison control weighted level integrals and signed endpoint terms.

If the normalized scalar integral were unbounded, rescale a
counterexample sequence and take a lower-curvature pointed limit.
Finite angular packing leaves finitely many exceptional points on
each compact set. Away from them, regular slabs and the
lower-dimensional estimate bound scalar integrals on annuli.
Concentrating at an exceptional point strictly increases the bounded
angular packing rank, so iteration is impossible. In codimension zero
the weighted estimate controls
\[
 \frac{\int_{B(p,1)}R^+\,d\mu}{1+\mu(B(p,2))};
\]
Bishop comparison bounds the denominator uniformly.

The three-dimensional construction starts from compact surface
levels. In dimension two, adjoin a flat line and compare the product
unit ball with $B(p,1/2)\times(-1/2,1/2)$. The bound $R\geq-2$
converts the result to an estimate for $R^+$; a finite half-ball
cover recovers the unit ball. Curvature is zero in dimensions zero
and one. These are the low-dimensional branches in the Lean
producer, followed by weighted corner induction in higher dimensions.
\end{proof}
The regular-slab and concentration arguments correspond to
\cite[Sections 3.1--3.7, pp.~5--7; Sections 4.3--4.7,
pp.~8--14]{Petrunin2009}.

\subsection{From volume to curvature}

\begin{lemma}[Unit-ball volume loss]\label{lem:rf-volume-loss}
\lean{PoincareMT.RicciFlow.exists_uniform_unitBall_volume_lower_bound}
\uses{thm:rf-scalar-integral,prop:rf-evolution}
There is a dimensional constant $P_n>0$ such that a Ricci flow
with complete slices and nonnegative curvature operator satisfies
\[
 \mu_{g(b)}(B_{g(b)}(x,1))
 \geq\mu_{g(a)}(B_{g(a)}(x,1))-P_n(b-a)
\]
whenever $[a,b]$ lies in the interior of its time domain.
\end{lemma}
\begin{proof}
Distances decrease, so the initial unit ball stays inside the evolving
unit ball. On the fixed initial ball,
$\partial_t d\mu=-R\,d\mu$. Its compact closure permits differentiation
under the integral. Since $R\geq0$, Theorem~\ref{thm:rf-scalar-integral}
bounds the loss rate by $P_n$. Integrate and use the terminal inclusion.
\end{proof}

\begin{lemma}[Reciprocal-time scalar bound]\label{lem:rf-time-scalar}
\lean{PoincareMT.RicciFlow.exists_time_mul_scalarCurvature_bound_of_unit_ball_volume}
\uses{lem:rf-volume-loss,thm:rf-compactness,prop:bounded-ancient-zero-volume}
Let $n\geq2$ and let $[a,b]$, with $a<b$, lie in the interior
of a complete Ricci flow with nonnegative curvature operator.
If $\mu_{g(t)}(B_{g(t)}(x,1))\geq\nu>0$ for every $x$ and
$t\in[a,b]$, then there is $K$ such that
\[
                (t-a)R(x,t)\leq K\qquad(a<t\leq b,\ x\in M).
\]
\end{lemma}
\begin{proof}
Otherwise point picking and parabolic rescaling give expanding
controlled cylinders with a backward time buffer. The local
derivative estimates and pointed compactness yield a complete,
nonflat ancient limit with bounded curvature and nonnegative
curvature operator. Volume comparison transfers the unit-ball
lower bound to every limit radius, with a positive constant
$\nu/4^n$. The limit consequently has positive asymptotic volume
ratio, contrary to Proposition~\ref{prop:bounded-ancient-zero-volume}.
Only the bounded ancient volume theorem is used here; its Harnack
input is Proposition~\ref{prop:rf-bounded-harnack}.
\end{proof}

\begin{lemma}[Local scalar estimate]\label{lem:rf-local-scalar}
\lean{PoincareMT.RicciFlow.exists_local_scalar_curvature_constant}
\uses{prop:rf-evolution,thm:rf-shi}
For $n\geq2$ there is $C_n>0$ with the following property.
Let a connected complete Ricci flow have nonnegative Ricci
curvature on $[a,b]$ inside its time domain. Let $r>0$ and $K\geq1$.
Assume $R(x,a)\leq r^{-2}$ on $B_{g(a)}(p,r)$ and
$R(x,t)\leq K/(t-a)$ on $B_{g(t)}(p,r)$ for $a<t\leq b$.
Then
\begin{equation}\label{eq:rf-chen-local}
 R(x,t)\leq
 \frac{\exp(C_nK)}{(r-d_{g(t)}(p,x))^2}
 \quad(t\in[a,b],\ x\in B_{g(t)}(p,r)).
\end{equation}
\end{lemma}
\begin{proof}
Apply the scalar evolution inequality to a distance cutoff on nested
balls. Shortened minimizing segments supply upper supports at the
cut locus. The reciprocal-time bound controls the time-dependent
cutoff errors. Choosing its exponential weight with coefficient
$C_nK$ absorbs those errors at a first positive maximum and yields
\eqref{eq:rf-chen-local}. The initial term is controlled by $r^{-2}$.
\end{proof}

\begin{lemma}[A short right-hand curvature bound]\label{lem:rf-right-bound}
\lean{PoincareMT.RicciFlow.exists_right_curvatureTensorNorm_bound_on_component_of_bounded_ancient_zero_volume}
\uses{lem:rf-initial-volume,lem:rf-volume-loss,lem:rf-time-scalar,lem:rf-local-scalar}
Suppose a complete Ricci flow has nonnegative curvature operator
on a component throughout $[a,b]$, with $a<b$ inside its time domain,
and bounded curvature operator there at time $a$. There are
$d\in(a,b]$ and $B\geq0$ with
$|\operatorname{Rm}|(x,t)\leq B$ on that component for $a\leq t\leq d$.
\end{lemma}
\begin{proof}
Lemma~\ref{lem:rf-initial-volume} gives $v>0$ initially.
Lemma~\ref{lem:rf-volume-loss} retains $v/2$ on a short right interval.
Lemma~\ref{lem:rf-time-scalar} gives the reciprocal-time scalar bound.
Apply Lemma~\ref{lem:rf-local-scalar} at every center, choosing one
radius from the initial uniform curvature bound. Nonnegative
curvature operator converts the scalar bound to a full-curvature bound.
In dimensions zero and one the curvature vanishes directly.
\end{proof}

\begin{lemma}[From local right bounds to a slab bound]
\label{lem:rf-slab-bound}
\lean{Poincare.RicciFlow.Harnack.exists_uniform_bound_of_local_right_bound_and_harnack}
\uses{prop:rf-bounded-harnack,lem:rf-right-bound}
Let $f:X\times[a,b]\to[0,\infty)$ be continuous in time at each
point and bounded on each slice. Suppose every $c<b$ has a
right neighborhood with one uniform bound, and on each bounded
initial slab $(t-a)f(x,t)$ is nondecreasing in $t$ for every $x$.
Then $f$ is uniformly bounded on $X\times[a,b]$.
\end{lemma}
\begin{proof}
Let $b_*$ be the supremum of terminal times of bounded initial
slabs, and choose one initial bounded slab $[a,c_0]$.
Monotonicity and continuity give
$(t-a)f(x,t)\leq(b_*-a)f(x,b_*)$ for $a<t\leq b_*$.
The slice bound at $b_*$ therefore bounds $f$ uniformly on
$[c_0,b_*]$. Together with the initial bound this includes $b_*$
in the bounded interval. Its local right bound excludes $b_*<b$.
\end{proof}

Apply this lemma to scalar curvature on each component.
Proposition~\ref{prop:rf-bounded-harnack} supplies the required
monotonicity on bounded slabs. Nonnegative curvature operator
again turns scalar control into full-curvature control. The
bounded-slab Hamilton argument now applies on every compact time
interval. Sending its left endpoint to $T_0$ proves the
differential assertion of Theorem~\ref{thm:rf-harnack}.
This is the volume-to-curvature passage of
\cite[Theorem 4.2 and Corollary 4.4, pp.~3160--3161]{CabezasRivasWilking2015}.

\subsection{Path integration and ancient times}

Assuming the differential estimate, apply it with $v=\gamma'/2$.
Nonnegative Ricci curvature implies
$\operatorname{Ric}(\gamma',\gamma')\leq R|\gamma'|^2$, so
for $f(s)=R(\gamma(s),s)$,
\[
 f'(s)\geq-\left(\frac1{s-T_0}
                       +\tfrac12|\gamma'(s)|^2_{g(s)}\right)f(s).
\]
Multiplication by
$(s-T_0)\exp(\tfrac12\int_{t_1}^s|\gamma'|^2)$ and integration proves
\eqref{eq:rf-harnack-path}.  This proof never divides by $R$, so it
also covers its zero set.  The finite-time factors must be retained.

\begin{corollary}[Ancient Harnack inequalities]\label{cor:rf-ancient-harnack}
\lean{Poincare.Geometry.RicciFlow.Harnack.ancient_differential_of_finite,Poincare.Geometry.RicciFlow.Harnack.ancient_integrated_of_finite_integrated}
\uses{thm:rf-harnack,prop:rf-evolution}
Let a Ricci flow on a connected $n$-manifold be defined for $t\leq0$.
Assume complete slices, nonnegative curvature operator, slice-wise
bounded curvature operator, and nonflatness on every slice
($\forall t\leq0\ \exists x:\ |\operatorname{Rm}|(x,t)\ne0$).
Then
\[
 \partial_tR+2dR(v)+2\operatorname{Ric}(v,v)\geq0
 \quad(t\leq0),
\]
with the derivative at zero taken from the left.  For every
$t_1<t_2\leq0$ and $C^1$ path,
\[
 R(\gamma(t_2),t_2)\geq R(\gamma(t_1),t_1)e^{-E(\gamma)/2}.
\]
In particular scalar curvature at a fixed spatial point is
nondecreasing in time.
\end{corollary}
\begin{proof}
For $t<0$ restrict to $(T_0,0)$ and let $T_0\to-\infty$ in
\eqref{eq:rf-harnack}; $R/(t-T_0)$ tends to zero and the derivative is
the same actual scalar derivative for all these restrictions.
Smoothness up to zero passes the inequality to the left derivative
there.  The previous integrating factor without the initial-time
factor proves the path estimate, including its endpoint by continuity.
Take $v=0$ for monotonicity.  These deductions also apply directly to
uniformly bounded ancient flows using only the bounded-slab theorem,
as required in the preceding reduction.
\end{proof}
Compare \cite[Theorem 4.40, p.~82]{MT2007}.

\section{Constructing harmonic coordinates}
\label{sec:rf-harmonic-coordinates}

Curvature bounds do not directly bound derivatives of the metric in
an arbitrary coordinate system.  Harmonic coordinates provide a
controlled choice.  We need the construction both for local geometric
analysis and, in dimension two, for the annular conformal coordinates
in Section~\ref{sec:annular-uniformization}.  The construction begins
on a radial coordinate ball; no global completeness assumption enters.
In this section $B_s$ denotes the Euclidean ball of radius $s$
centered at zero; other centers are displayed explicitly.

\begin{theorem}[Uniform harmonic coordinates]\label{thm:rf-harmonic-radius}
\lean{PoincareMT.RiemannianMetric.exists_uniform_harmonic_radius}
\uses{def:foundations-manifold}
Fix $n\geq2$, $R>0$, and $K\geq0$.  There are $r>0$ and $H>0$,
depending only on these data, with the following property.  Suppose
$g$ is a smooth positive metric on $\mathbb R^n$ satisfying
\[
 \tfrac14|v|^2\leq g_x(v,v)\leq\tfrac94|v|^2,
 \qquad |\operatorname{Rm}_g|(x)\leq K\quad(x\in B_R),
 \qquad g_x(x,w)=\langle x,w\rangle\quad(x,w\in\mathbb R^n).
\]
There is a smooth local diffeomorphism $e:B_{2r}\to B_R$ with
$e(0)=0$ such that $h=e^*g$ extends to a smooth positive metric on
$\mathbb R^n$.  On $B_{2r}$ the standard coordinate functions
satisfy $\Delta_h y_i=0$, and
\[
 \tfrac19 I\leq h\leq9I,\qquad
 |Dh(x)|\leq H,\qquad
 |Dh(x)-Dh(y)|\leq H|x-y|^{1/2}.
\]
Moreover $d_g(0,e(x))\leq3|x|$.  The derivatives and norms of
the coefficient field $h$ in this statement are Euclidean ones.
\end{theorem}

The comparison with the usual harmonic-coordinate theorem is local:
\cite[Proposition 26.49, Step 3, equations (26.146)--(26.150),
pp.~383--384]{ChowEtAl2010} obtains such coordinates on an exponential
lift by invoking Jost--Karcher and DeTurck--Kazdan.  We develop the
radial-metric construction used here.  Its constants precede the
choice of $g$; qualitative smoothness constants used along the way
need not be uniform.

\subsection{Radial transport and weak harmonic replacements}
\label{sec:rf-weak-harmonic-replacements}

Write $T(x)v$ for parallel transport of $v$ along $t\mapsto tx$,
$0\leq t\leq1$.  The Gauss identity implies $g_0=I$ and makes
these radial lines affinely parametrized geodesics.  Thus $T(x)$
is an isometry from the Euclidean inner product to $g_x$.  Varying
the radial line in direction $w$ and commuting the two covariant
derivatives gives
\[
 T(x)^{-1}\nabla_w(T(\cdot)v)|_x
 =\int_0^1T(tx)^{-1}
       \operatorname{Rm}_{tx}(x,tw)T(tx)v\,dt.
\]
Here the curvature is regarded as its $(1,3)$ tensor with the
same convention as the connection commutator.  The metric upper
bound gives, with $A=3/2$,
\begin{equation}\label{eq:rf-radial-frame}
 |\nabla_w(T(\cdot)v)|_g\leq KA^2|x||w||v|,
 \qquad |T(x)^{-1}w-w|\leq KA^2|x|^2|w|.
\end{equation}
For the second inequality, differentiate the radial coframe and
use torsion freeness to interchange the radial and transverse
directions.  Its difference from the identity is the integral of
the frame connection evaluated on $(w,tx)$.  The first bound,
integrated along the same segment, gives the second.  Consequently
$g(v,v)=|T^{-1}v|^2$ differs from $|v|^2$ by at most
$\theta(2+\theta)|v|^2$ wherever $KA^2|x|^2\leq\theta$.

Put $a(x)=\sqrt{\det g(x)}\,g(x)^{-1}$.  On a sufficiently
small ball $B_{4s}$ the preceding estimate makes $a-I$ smaller
than any prescribed $\varepsilon>0$.  This choice depends only
on $n,R,K,\varepsilon$.  On $B_{2s}$ seek
$U_i=x_i+w_i$, with $w_i\in H^1_0(B_{2s})$, satisfying
\begin{equation}\label{eq:rf-harmonic-replacement}
 \int a\nabla w_i\cdot\nabla\phi
   =-\int(a-I)e_i\cdot\nabla\phi
       \quad(\phi\in H^1_0(B_{2s})).
\end{equation}
All integrals here are with respect to Euclidean measure.  The
subtracted identity contributes zero because $\phi$ has zero
boundary trace.  Ellipticity and the ball Poincare inequality
make the left side a coercive bounded form on the zero-boundary
energy completion.  The Riesz representation theorem therefore
constructs a unique $w_i$.  Testing with $w_i$ and then applying
Poincare gives
\begin{equation}\label{eq:rf-correction-energy}
 \int_{B_{2s}}|\nabla w_i|^2\leq C\varepsilon^2s^n,
 \qquad \int_{B_{2s}}|w_i|^2\leq C\varepsilon^2s^{n+2}.
\end{equation}
Equivalently, $U_i$ is weakly harmonic for $g$.  In constructing
this completion with the Riemannian volume, the linear boundary
value is first multiplied by a compact cutoff equal to one on
the ball.  Locality removes this cutoff from the equation; metric
and volume comparability identify the displayed energy norms.

Each $U_i$ has a smooth interior representative before any
quantitative coordinate estimate is applied.  Here is the local
regularity mechanism.  The derivatives in the energy completion
are weak coordinate derivatives, as follows by integration by
parts on a sequence of smooth approximants.  In a relatively
compact subball, test the divergence equation by the backward
difference quotient of $\eta^2$ times a forward difference
quotient of $U_i$.  Coercivity bounds the first derivatives of
these difference quotients; derivatives of $a$ and $\eta$ occur
only in bounded lower-order terms.  Weak compactness yields
second weak derivatives.  Differentiating the distributional
equation repeats this argument on smaller balls, giving every
local Sobolev order.  Sobolev embedding then gives a smooth
representative.  Two such representatives agree everywhere on
an overlap because they agree almost everywhere and are
continuous; positive volume of open sets proves this last step.
The original weak equation consequently becomes $\Delta_gU_i=0$
pointwise.  This argument uses bounded derivatives of this
particular smooth metric, and asserts no uniform estimates yet.

\subsection{From small energy to an invertible coordinate map}

The energy bound alone does not make $DU$ close to the identity.
To obtain that conclusion, compare $\nabla U_i$ with the radial
unit field $X_i=T e_i$, and write $V=\nabla U_i-X_i$.
For any compactly supported vector field $Y$, commuting derivatives
of the harmonic scalar gives the weak connection identity
\[
 \int\langle\nabla V,\nabla Y\rangle_g
 =-\int\operatorname{Ric}(\nabla U_i,Y)
   -\int\langle\nabla X_i,\nabla Y\rangle_g.
\]
The sign of the Ricci term follows from
$\Delta^\nabla\nabla U_i=\operatorname{Ric}^{\sharp}\nabla U_i$.
Choose $\sigma=C(n)(K+1)s^2$ so that
$s^2 nK\,|X_i|_g\leq\sigma$ and
$|\nabla X_i|_g\leq\sigma/s$, using
\eqref{eq:rf-radial-frame}.  Test the identity with
$\eta^2(|V|_g^2+\sigma^2)^{p-1}V$, $p\geq1$.
The derivative of the power has the favorable sign; Cauchy's
inequality absorbs half the derivative energy.  With
$v=(|V|_g^2+\sigma^2)^{1/2}$ this gives
\begin{equation}\label{eq:rf-power-energy}
 \int|\nabla(\eta v^p)|_g^2
 \leq22p^2\left(
    \frac{nK+2}{s^2}\int\eta^2v^{2p}
       +\int v^{2p}|\nabla\eta|_g^2\right).
\end{equation}
These integrals use $d\mu_g$ and $s\leq1$.

For clarity, the iteration retains the radius dependence.
Ellipticity transfers the Euclidean Sobolev inequality with
exponent $2\chi=2n/(n-1)$, which also applies when $n=2$.
Apply it to the cutoffs in \eqref{eq:rf-power-energy}, with
$p=\chi^j$ and balls about $z$ of radii $s/2^{j+1}$.
If $b_j$ is the $L^{2\chi^j}$ norm on the $j$th ball, then
\[
 b_{j+1}\leq((B/s)L^j)^{1/(2\chi^j)}b_j.
\]
The sums of $\chi^{-j}$ and $j\chi^{-j}$ converge, and
$\sum_j(2\chi^j)^{-1}=n/2$.  The product is therefore bounded
by $Cs^{-n/2}$.  Continuity at $z$, together with the shrinking
balls and growing exponents, gives
\[
 |V(z)|_g^2+\sigma^2
 \leq Cs^{-n}\left(
       \int_{B(z,s/2)}|V|_g^2\,dx
           +|B(z,s/2)|\sigma^2\right)
       \quad(z\in B_{s/2}).
\]
The measure in this last formula can be Euclidean by the fixed
volume bounds.  By \eqref{eq:rf-correction-energy} and the coframe
estimate, its right side is at most
$C\varepsilon^2+C(K+1)^2s^4$.
The same coframe estimate compares $X_i$ with $\nabla x_i$;
ellipticity converts the gradient error to an ordinary differential
error.  Thus, given any $\tau>0$, choosing $\varepsilon$ and then
$s$ sufficiently small yields
\[
 |DU_i-e_i^*|\leq\tau\quad\hbox{on }B_{s/2}.
\]
All these choices are made before the metric.

Center the map $F=(U_1,\ldots,U_n)-U(0)$.  Choose the component
tolerance so that $|DF-I|\leq\beta\leq1/2$ on a ball $B_\rho$.
Integrating along segments shows that $F-I$ is
$\beta$-Lipschitz.  Hence $F$ is injective and its differential
is invertible.  For $|y|<\rho/4$, the contraction
$x\mapsto y-(F(x)-x)$ preserves $\overline B_{\rho/2}$ and
has a unique fixed point.  This proves
$B_{\rho/4}\subset F(B_\rho)$, with a smooth inverse on the
image.  The inverse differential has bounds $2/3$ and $2$.
Its pullback metric therefore has ellipticity $1/9,9$.
More precisely, if $|g-I|\leq\alpha$, expansion of
$(DF^{-1})^Tg\,DF^{-1}$ gives
\[
 |(F^{-1})^*g-I|\leq4\alpha+6\beta.
\]
Both errors can be chosen arbitrarily small.  This is the
quantitative freedom needed for the next estimate.

\subsection{Uniform first derivatives without curvature derivatives}

Now let $h$ be a harmonic-coordinate metric on $B_S$, with
$aI\leq h\leq bI$ and $|\operatorname{Rm}_h|\leq K$.
For a harmonic function $f$ with $|\nabla f|_h\leq G$, put
$Q=\nabla^2f$.  The differentiated equation is used in the form
\begin{equation}\label{eq:rf-hessian-divergence}
 \Delta^\nabla Q=\operatorname{div}\mathcal B+\mathcal R(Q),
 \qquad
 \mathcal B(a,b,c)=h(a,b)\operatorname{Ric}(\nabla f,c)
                         -df(\operatorname{Rm}(a,b)c).
\end{equation}
Here $\mathcal R(Q)$ is the algebraic curvature action from
commuting a traced derivative with a free tensor slot:
\[
 \mathcal R(Q)(a,b)=-\sum_j\left[
 Q(\operatorname{Rm}(E_j,a)E_j,b)
       +Q(E_j,\operatorname{Rm}(E_j,a)b)\right]
\]
for an orthonormal frame $E_j$.  The identity follows by writing
the Codazzi defect of $Q$ as $-df(\operatorname{Rm}(a,b)c)$
and the divergence of $Q$ as $\operatorname{Ric}(\nabla f,\cdot)$.
In particular,
$|\mathcal B|\leq C(n)KG$ and
$|\mathcal R(Q)|\leq C(n)K|Q|$.
Integration by parts leaves the derivative on the tensor test:
\[
 \int\langle\nabla Q,\nabla Z\rangle
 =\int\langle\mathcal B,\nabla Z\rangle
       -\int\langle\mathcal R(Q),Z\rangle.
\]
This is why no bound for $\nabla\operatorname{Rm}$ is required.

Set $v=(|Q|^2+\sigma^2)^{1/2}$, with
$\sigma=(1+C(n)KG)s$ on a ball of radius $s\leq1$.
Testing by $\eta^2v^{2p-2}Q$ gives the same power estimate as
\eqref{eq:rf-power-energy}, with $nK+2$ replaced by
$2n^2K+1$.  The preceding iteration bounds $|Q(z)|^2$ by
its local squared mass and $\sigma^2$.  To supply that mass,
integrate the scalar Bochner identity
\[
 \tfrac12\Delta|\nabla f|^2
       =|Q|^2+\operatorname{Ric}(\nabla f,\nabla f)
\]
against a squared cutoff.  Absorbing the mixed cutoff term yields
\[
 \int_{B_{S/4}}|Q|^2
     \leq C(a,b,n)(K+S^{-2})G^2|B_{S/2}|.
\]
The mass and iteration bounds give a uniform Hessian bound on
$B_{S/8}$.  For $f=y_i$, the coordinate gradient bound is
$G\leq a^{-1/2}$ and
$\nabla^2 y_i(\partial_j,\partial_k)=-\Gamma^i_{jk}$.
Thus the connection coefficients are bounded.  Metric
compatibility,
\[
 \partial_k h_{ij}
       =h_{\ell j}\Gamma^\ell_{ki}
                          +h_{i\ell}\Gamma^\ell_{kj},
\]
gives $|Dh|\leq L(n,S,a,b,K)$ on this smaller ball.

\subsection{The half-Holder estimate and the local surface chart}

Harmonicity of the coordinates cancels the contracted Christoffel
terms in the Ricci equation.  For each coefficient $u=h_{ij}$,
\[
 \Delta_h u=-2\operatorname{Ric}_{ij}+\mathcal Q_{ij}(Dh,Dh),
 \qquad
 |\Delta_hu|\leq2nKb+\frac{32n^2b}{a^3}L^2.
\]
The last bound follows by expanding the quadratic connection
terms and using the inverse-metric bound $a^{-1}$.  Put
$A=\sqrt{\det h}\,h^{-1}=I+E$.
The map $h\mapsto\sqrt{\det h}\,h^{-1}$ is smooth near $I$,
so the freely chosen bound $|h-I|<\delta$ makes $|E|$ as small
as desired, while the established bound on $Dh$ bounds $DE$.
Choose a spatial cutoff $\eta$ supported in $B_{S/32}$ and
equal to one on $B_{S/128}$.  For $v=\eta u$ the equation is
\[
 \Delta v=F-\operatorname{div}(E\nabla v),
\]
where $F=\operatorname{div}(A\nabla v)$ has a uniform supremum
bound.  Indeed its expansion contains only the bounded
$\sqrt{\det h}\Delta_hu$, $u$, $Du$, $A$, $DA$, and the first
two cutoff derivatives.  A second cutoff extends $E$ globally,
preserving its small supremum and a bounded half-Holder seminorm.

Here is the estimate that absorbs the last term.  Multiply $v$
by a fixed smooth time cutoff $\zeta$, zero at time zero and
one at time one.  The resulting function has heat residual
\[
 (\partial_t-\Delta)(\zeta v)
    =\zeta'v-\zeta F+
                        \operatorname{div}(\zeta E\nabla v).
\]
The zero-initial-data heat representation bounds its gradient
half-Holder seminorm at time one by
\[
 C\bigl(\|v\|_\infty+\|F\|_\infty
                    +n[E\nabla v]_{1/2}\bigr).
\]
For the divergence term this follows from cancellation of the
second derivative of the Gaussian kernel.  Subtract the source
value at the center in the convolution.  If its half-Holder
seminorm is $Q_0$ and $d=|x-y|$, the difference of the two
Hessian convolutions is bounded both by
$C Q_0t^{-3/4}$ and by $C Q_0d\,t^{-5/4}$; the latter uses
the third derivative of the kernel.  Splitting the time
integral at $d^2$ gives $C Q_0d^{1/2}$.  The ordinary bounded
source is estimated using the first kernel derivative.
Consequently, with $J=[Dv]_{1/2}$,
\[
 J\leq C(\|v\|_\infty+\|F\|_\infty
                       +n[E]_{1/2}\|Dv\|_\infty)
                           +Cn\|E\|_\infty J.
\]
Compact smoothness makes $J$ finite initially.  Choosing the
coefficient tolerance so that $Cn\|E\|_\infty\leq1/2$
absorbs it, giving a uniform bound independent of this initial
seminorm.  Summing the finitely many coefficient slots proves
the asserted bound on $[Dh]_{1/2}$.

To finish Theorem~\ref{thm:rf-harmonic-radius}, choose this small
tolerance first, then the weak-coordinate radii and inverse map.
Take $S=\rho/8$ in its image and $r=S/256$.
Extend $(F^{-1})^*g$ by a convex cutoff combination with the
Euclidean metric, equal to the pullback on $B_S$ and with
ellipticity $1/9,9$ globally.  On $B_S$, change of coordinates
preserves curvature and identifies $\Delta_h y_i$ with
$\Delta_gF_i=0$.  The preceding estimates apply on $B_{S/128}
=B_{2r}$.  Finally the image under $e=F^{-1}$ of the segment
$tx$ has length at most $3|x|$, proving the distance bound.

For an arbitrary smooth surface metric near $p$, choose a
normalized exponential chart.  Shrink it until its coefficients
are within $1/2$ of $I$, and extend them by a convex cutoff
with $I$.  The Gauss identity survives this extension.  On a
smaller compact ball the smooth curvature has a finite bound,
so the theorem supplies a harmonic chart centered at $p$.
Its first coordinate $u$ has nonzero differential.  On a small
oriented disk the one-form $\iota_{\nabla u}d\mu$ is closed,
because $\Delta u=0$, and hence equals $dv$ for a smooth
function $v$.  The covectors $du,dv$ are orthogonal, have equal
nonzero norm, and give an invertible differential of $(u,v)$.
The inverse function theorem therefore gives coordinates in
which the metric is a positive smooth multiple of $du^2+dv^2$.
This proves the local isothermal input used for the annulus,
including the absence of any completeness premise on that
surface metric.

\section{Pointed compactness and complete interior slices}
\label{sec:rf-compactness}

\begin{definition}[Pointed smooth convergence]\label{def:rf-pointed-convergence}
\lean{PoincareMT.PointedGeometricConvergence}
\uses{def:rf-flow}
A sequence of based flows $(M_k,g_k(t),p_k)$ on $J=(T',T)$ converges
smoothly in the pointed sense after a subsequence if there are a
connected based flow $(M_\infty,g_\infty(t),p_\infty)$, a strictly
increasing subsequence $\sigma$, connected open sets
$U_j\subset U_{j+1}$ with compact closures and union $M_\infty$,
and base-preserving smooth open embeddings
$\phi_j:U_j\to M_{\sigma(j)}$, such that
$\phi_j^*g_{\sigma(j)}\to g_\infty$ with every mixed space and time
derivative uniformly on compact subsets of $J\times M_\infty$.
Each compact subset is required to lie in the domains for all
sufficiently large $j$.  The equivalent spacetime maps
$(t,x)\mapsto(t,\phi_j(x))$ preserve time and its vector field
$\partial_t$.
\end{definition}

\begin{theorem}[Pointed Ricci-flow compactness]\label{thm:rf-compactness}
\lean{PoincareMT.pointedRicciFlowCompactness}
\uses{def:rf-pointed-convergence,thm:rf-shi,prop:rf-metric-comparison}
Fix $n$ and $T'<0<T$.  Let $(M_k,g_k(t),p_k)$ be connected based
$n$-dimensional Ricci flows on $J=(T',T)$, with calibrated Riemannian
volume on the zero slice.  Suppose:
\begin{enumerate}
 \item For every $A>0$, $\overline{B_{g_k(0)}(p_k,A)}$ is compact
 for all sufficiently large $k$.
 \item For every $A>0$ there is $K_A\geq0$ such that, for all
 sufficiently large $k$, for every $s,t\in J$ and every
 $x\in B_{g_k(s)}(p_k,A)$,
 $|\operatorname{Rm}_{g_k(t)}|(x)\leq K_A$.
 \item There are $r_0,\kappa>0$ such that, eventually,
 $\mu_{g_k(0)}(B_{g_k(0)}(p_k,r_0))\geq\kappa r_0^n$.
\end{enumerate}
There is a smoothly pointed convergent subsequence, and its limit is
complete at every $t\in J$.
\end{theorem}

The spacetime formulation additionally records, on each zero-time ball
and each compact time interval containing zero, a time-preserving
embedding carrying $\partial_t$ to $\partial_t$, equal to the identity
at time zero, with a curvature bound on its image.  For these ordinary
product flows, the identity map supplies this datum from condition~2.
No completeness of an individual approximating manifold is assumed.
The eventual index may depend on $A$; $K_A$ is independent of both
times and of the eventual index.

\subsection{Charts, overlaps, and the limit metric}

We describe the geometric construction rather than presuppose the
limit manifold.  The case $n=0$ is the one-point connected manifold.
For $n>0$, fix a radius and a compact time interval in $J$.
Local curvature bounds and the volume lower bound give a lower
injectivity radius on a slightly smaller zero-time ball.
Volume comparison propagates the base-ball lower bound to centers
at bounded distance.  Normal coordinates then have positive lower
and upper metric bounds.  Derivative estimates on a buffered time
interval, metric comparison, and differentiated coordinate equations
bound every mixed jet of their metric coefficients.  The radius of
each chart and every constant may depend on the chosen ball and
buffered interval.

We spell out the uniform size of the chart cover.  For a fixed source
metric write $B(p,A)$ for the ball to be covered, and fix
$0<\delta\leq A$.  Suppose $\overline{B(p,5A)}$ is compact and
$\operatorname{Ric}\geq-(n-1)Kg$ there, with $K\geq0$.  Let
\[
 V_{n,K}(r)=|S^{n-1}|\int_0^r s_K(u)^{n-1}\,du,\qquad
 s_K(u)=
 \begin{cases}
 u,&K=0,\\
 \sinh(\sqrt K u)/\sqrt K,&K>0.
 \end{cases}
\]
Thus $V_{n,K}$ is the ball-volume function for the simply connected
constant-curvature $-K$ model, with $V_{1,K}(r)=2r$.
For each $x\in B(p,A)$, the ball $B(x,4A)$ has compact closure
inside $\overline{B(p,5A)}$, and $B(p,2A)\subset B(x,3A)$.
Local relative volume comparison gives
\[
 \mu(B(x,\delta/2))
 \geq\frac{V_{n,K}(\delta/2)}{V_{n,K}(3A)}
                          \mu(B(x,3A))
 \geq\frac{V_{n,K}(\delta/2)}{V_{n,K}(3A)}
                          \mu(B(p,2A)).
\]
The local form of Bishop--Gromov uses compact closure of the larger
ball, so its radial geodesics and polar integration stay in the
controlled region; completeness of the whole manifold is unnecessary
(\cite[Theorem 1.34, p.~19; proof of Theorem 5.6, p.~86]{MT2007}).

If a finite set of centers in $B(p,A)$ has pairwise distance at least
$\delta$, its open $\delta/2$-balls are disjoint and lie in $B(p,2A)$.
The latter has positive finite volume.  Adding the displayed lower
bounds and cancelling that volume shows that the number of centers
is at most
\[
                         \frac{V_{n,K}(3A)}{V_{n,K}(\delta/2)}.
\]
Start with any center and add a point outside all open $\delta$-balls
whenever such a point exists.  Every addition preserves the separation,
so this process must terminate before exceeding the displayed bound.
The terminal $\delta$-balls cover $B(p,A)$.  For $\delta>A$ the
single basepoint ball already covers it.
Condition~1 and the zero-time case of condition~2 make the bound
uniform along the eventual source sequence.  This packing step does
not use the noncollapse hypothesis; noncollapse is needed for the
uniform injective coordinate radius.

Here is the noncollapse-to-injectivity step used for that radius.
Relative volume comparison first transports
$\mu(B(p,r_0))\geq\kappa r_0^n$ to every center in $B(p,A)$:
choose a larger controlled ball, compare the center's small ball to a
fixed base ball, and retain a positive lower bound depending only on
$A,r_0,\kappa$ and the curvature bound.  On the same controlled ball
the curvature tensor has a uniform bound $K$.  The Jacobi equation
then gives a common radius $s$ below both the curvature comparison
radius and the controlled source radius.  The exponential map at each
center is nonsingular on $B(0,s)$.

Injectivity is a separate global issue.  The local metric argument
in the proof of Lemma~\ref{lem:ancient-positive-ratio-obstruction},
starting at \eqref{eq:ancient-local-injectivity}, applies here with
the transported volume lower bound; that argument uses only local
curvature, volume and compact confinement.  In detail, relative volume
comparison gives a lower bound $w_0>0$ for the center's $s/4$-ball,
whereas the Jacobi bound gives an upper bound
$P=n!(3/2)^n\omega_ns^n$ for the pullback volume of $B(0,s)$.
Choose $N$ with $(N+1)w_0>P$, then choose the possible collision
radius $\rho$ small enough for the short-return construction with
this $N$.  A first collision gives a radial return of length at most
$2\rho$.  Its successive lifts have $N+1$ distinct endpoints:
a repeated endpoint would give a finite orbit of a local isometry,
whose sum of squared radial norms has two minima along a strictly
convex geodesic.  The explicit radius margins and this convexity
argument are given there.  Lifting radial paths from these endpoints
gives at least $N+1$ preimages of every point of the $s/4$-ball.
Change of variables would then imply $(N+1)w_0\leq P$, a
contradiction.  Thus the closed $\rho$-ball is injective.
Restricting the canonical exponential to the open
ball makes it a diffeomorphism onto the intrinsic ball of radius
$\rho$; radial length equality and the normal-coordinate derivative
bound give the two-sided distance estimates used above.

Choose a normal-coordinate radius $R$ uniformly for all centers in
$B(p,A)$, and a smaller $\rho>0$ with $2\rho<R$ on which the metric
coefficient estimates hold.  Apply the preceding cover with
$\delta=\rho/4$.  Radial distance in an injective normal chart is
exactly the Euclidean norm, so its $\rho/4$-ball is covered by the
image of the compact coordinate ball $\overline{B(0,\rho/4)}$.
Insert the basepoint among the centers, increasing the bound by at
most one, and repeat centers as necessary to label every source cover
by the same finite set with index zero at the basepoint.  The maps
and centers vary with the source index, but the coordinate domains,
compact covering pieces and number of labels do not.

The ambient distance bounds needed below also deserve attention.
Suppose $a|\xi|^2\leq\Phi^*g(\xi,\xi)\leq b|\xi|^2$ on
$B(0,2\rho)$, with $a>0$.  For $x,y\in B(0,\rho/2)$, the straight
coordinate segment gives
$d_g(\Phi(x),\Phi(y))\leq\sqrt b\,|x-y|$.
For the opposite inequality choose paths whose lengths approach this
ambient distance.  That distance is less than $\rho$, by joining the
two radial segments through the chart center, so the paths can be
chosen with length less than $\rho$.  Each stays in the intrinsic
$3\rho/2$-ball of the center and thus in the normal chart.
The inverse differential has norm at most $a^{-1/2}$ there.
Integration along the paths and passage to their infimum gives
\[
       \sqrt a\,|x-y|\leq d_g(\Phi(x),\Phi(y))
                         \leq\sqrt b\,|x-y|.
\]
Thus a metric lower bound in a larger coordinate region really does
control ambient distance on the inner region, even for incomplete
source metrics.

The coefficient compactness is joint in space and time.  For later use,
once the limit charts and embeddings below have been constructed,
take the limit chart matching source label $j$ and write its
coordinate parametrization as $q=q_j:V_j\to M_\infty$.
On compact subsets of $J\times V_j$, the pulled-back coefficients of
$g_k(t)$ have bounds for every mixed derivative.  The source readout
is first obtained locally: compact containment puts
$F_k(q_j(x))$ in the image of the matching source chart $e_{kj}$,
and its inverse
coordinate map is smooth there.  The readout converges to the identity
in every spatial jet.  Pulling back a bilinear coefficient field
therefore gives
\[
 \bigl((F_k\circ q)^*g_k(t)\bigr)_x(v,w)
 =G_{k,j}(t,b_k(x))
     \bigl(Db_k(x)v,Db_k(x)w\bigr),
\]
where $b_k\to\mathrm{id}$ smoothly and $G_{k,j}$ is the original
source-chart coefficient.  The chain rule and the product rule give
convergence of every mixed $(t,x)$ derivative.  A finite cover of a
compact set supplies one threshold for all local formulas and all basis
components; a diagonal choice over derivative order gives the
countable all-jets statement.  This is why the construction needs
spatial readout jets rather than only $C^0$ distance approximation.

For completeness, the bounds feeding this compactness have two
separate parts.  The local Shi estimate gives, for every curvature
derivative order $\ell$, a tail bound on a larger source ball and a
short time interval ending at the chosen interior time.  At zero time
we intersect the finitely many tails $\ell\leq m$ and use the
finite-order pullback estimate for a normal exponential chart:
curvature derivatives, the normalized initial metric and the radial
Gauss identity bound the $m$th metric-coefficient jet on the closed
coordinate ball.  The coupled radial-frame calculation leading to
\eqref{eq:ancient-radial-connection} gives this estimate: curvature
components, the Jacobi coframe and the connection are bounded in
that order, with spatial order $m$ of $\nabla^\ell\operatorname{Rm}$
requiring only $\ell+m$ covariant curvature derivatives.
This is a local calculation independent of the ancient-flow
application in which it is developed.  The threshold may depend on $m$, as it must for
$C^\infty$ convergence, but it is uniform over every center and every
normal chart in the fixed finite cover.

To extend the coefficient bounds through the whole time interval,
keep the spatial point and chart fixed.  If $w$ is the coordinate
tangent vector, the Ricci contraction satisfies
\[
 |\operatorname{Ric}_t(w,w)|
 \leq n^3K\,g_t(w,w)
\]
whenever $|\operatorname{Rm}_t|\leq K$.  The metric evolution
$\partial_tg_t=-2\operatorname{Ric}_t$ and Gronwall therefore give
\[
 e^{-2n^3K(T-T')}g_0(w,w)
 \leq g_t(w,w)
 \leq e^{2n^3K(T-T')}g_0(w,w).
\]
The curvature bound is required at the fixed image point for every
time; no uniform choice of a time slice is silently substituted.
Differentiating the metric equation and using the Shi bounds then
supplies the mixed time derivatives used by the spacetime jet
extraction.  This is the time-comparison mechanism in
\cite[Proposition 5.14, pp. 90--91; Theorem 5.15, pp. 91--92]{MT2007}.

Take a diagonal subsequence over increasing balls, compact time
intervals exhausting $J$, and derivative orders.  There are countably
many fixed coordinate domains $V_i$, source charts
$e_{ki}:V_i\to M_k$, smooth positive metric coefficient limits
$G_i(t,x)$, and limits
\[
                 D_{ij}(x,y)=\lim_k d_{g_k(0)}(e_{ki}(x),e_{kj}(y)).
\]
The distance convergence is locally uniform.  Uniform upper and lower
Lipschitz bounds for each chart pass to the limit; triangle symmetry
and the triangle inequality pass pointwise.  Identify $(i,x)$ with
$(j,y)$ exactly when $D_{ij}(x,y)=0$.  The triangle inequality proves
transitivity, and the lower chart bound shows that no two different
points of one chart are identified.

The zero-distance relation determines actual open overlaps.  Write
$c_i|x-x'|\leq D_{ii}(x,x')\leq L_i|x-x'|$ for the inherited
chart bounds.  First observe a source coverage fact.  If
$\overline B(y,r)\subset V_j$, the image of $B(y,r)$ under
$e_{kj}$ contains the source ball $B(e_{kj}(y),c_jr)$.
Indeed its image is open, its closure is compact, and any boundary
point is the image of a coordinate point at distance $r$ from $y$,
so lies at source distance at least $c_jr$.  Intersecting with the
connected source ball proves coverage.  These source balls are
connected because the Riemannian distance is a length distance;
near-minimizing paths to the center stay in the ball.  Completeness
is not needed.

Now suppose $D_{ij}(x,y)=0$.  Choose such a compact target ball and
a small coordinate neighborhood of $x$, of radius at most
$c_jr/[2(L_i+1)]$.  For all sufficiently large $k$, its whole
source image lies in $e_{kj}(B(y,r))$.  Thus the source inverse
transition is defined on this fixed neighborhood with compact target.
For any point $x'$ in the neighborhood, compactness of its inverse
coordinates gives a limiting partner $y'$ with $D_{ij}(x',y')=0$.
The partner is unique by the lower chart bound.  Hence the overlap
is open, and its transition $T_{ij}$ satisfies
\[
 |T_{ij}(x)-T_{ij}(x')|\leq (L_i/c_j)|x-x'|.
\]
Interchanging the charts gives its inverse; the triangle inequality
and uniqueness give every cocycle identity.  Compact subsets of an
overlap have one compact target and one eventual source threshold,
by a finite cover of these neighborhoods.

There is also quantitative smoothness.  A source transition $T_k$
is a local isometry between the coefficient metrics $G_{k,i}(0)$
and $G_{k,j}(0)$.  Uniform ellipticity and the upper metric bound
first bound $DT_k$.  Preservation of the Levi-Civita connection gives
\[
 D^2T_k=DT_k\,\Gamma_{k,i}
       -\Gamma_{k,j}(T_k)(DT_k,DT_k).
\]
Inverse metric bounds and the coefficient jets bound the Christoffel
symbols and their derivatives; differentiating this identity
inductively bounds every derivative of $T_k$ on the confined domain.
Smooth compactness on smaller compact sets now applies.  Every
subsequential limit must be $T_{ij}$ by the zero-distance identity,
so all transition jets converge to those of this one smooth map.

The quotient charts are open embeddings, since their overlap domains
are open.  The continuous functions $D_{ij}$ descend, by the triangle
inequality, to a separating distance on the quotient.  This distance
induces precisely the chart topology: distance balls are open in all
charts; conversely, a sufficiently small distance ball about $q_i(x)$
lies in the image of any chosen smaller compact coordinate ball about
$x$, by the same source coverage and compact-partner argument.
Thus the quotient is Hausdorff.  The countably many Euclidean chart
bases give a countable base, so the smooth transitions define a
smooth manifold.  This auxiliary distance is not yet identified with
the path distance of the metric constructed next.
The $G_i$ obey the tensor transition law, since it holds before taking
limits:
\[
 G_j(t,T_{ij}(x))(DT_{ij}(x)v,DT_{ij}(x)w)=G_i(t,x)(v,w).
\]
The compact target confinement just established permits evaluation
at the varying source transitions before passing to this identity.
The retained lower metric bounds make every $G_i$ positive, so these
coefficients give a smooth metric $g_\infty(t)$.  The coordinate
Ricci expression uses the metric inverse and at most two spatial
derivatives.  Positivity makes inversion continuous on each compact
coordinate neighborhood.  Convergence of two spatial derivatives and
one time derivative therefore passes
\eqref{eq:rf-equation} to the limit.  All mixed jets converge, so the
limiting flow is smooth.

Charts into each source do not initially agree exactly on the limiting
overlaps.  We construct actual embeddings near any given compact
quotient set by a finite induction.  Write $q_i:V_i\to M_\infty$
for the quotient charts.  The induction retains a compact set $P$,
an open neighborhood $O$ with compact closure, and embeddings
$F_k:O\to M_k$ for all sufficiently large $k$.  It also retains the
following local formula: near each point of $O$, in some fixed chart,
\[
             F_k(q_j(y))=e_{kj}(b_k(y)),
             \qquad b_k\longrightarrow\mathrm{id}\quad\hbox{in }C^\infty
\]
on compact sets.  A fixed compact neighborhood of the basepoint is
included in $P$ and its original source-chart formula is kept exactly.
One source chart initializes these properties.

These local formulas imply two useful convergence statements.
First, on every compact $C\subset V_i$ with $q_i(C)\subset O$,
\[
             d_{g_k(0)}(F_k(q_i(x)),e_{ki}(x))\longrightarrow0
             \quad\hbox{uniformly for }x\in C.
\]
Indeed, in a local formula choose $y$ with $q_j(y)=q_i(x)$.
The distance from $e_{kj}(b_k(y))$ to $e_{kj}(y)$ tends uniformly
to zero by the upper Lipschitz bound.  The distance from
$e_{kj}(y)$ to $e_{ki}(x)$ tends to $D_{ji}(y,x)=0$.
A finite cover gives the assertion on $C$.
Second, the readout $e_{ki}^{-1}F_kq_i$ converges smoothly to
the identity wherever $q_i$ takes values in $O$.  To verify that
the inverse is actually defined, take a compact coordinate ball
inside $V_i$.  The lower chart Lipschitz bound places its boundary
a fixed positive source distance from its center.  The image is
open, and its compact closure is closed; connectedness of small
source metric balls therefore puts a uniformly smaller source
ball inside that image.  The preceding distance approximation
places $F_kq_i$ in it.  This argument is local and does not require
complete source metrics.

For smooth convergence, let $T_{ij}$ be the limiting transition
from chart $i$ to chart $j$, and $T_{k,ji}=e_{ki}^{-1}e_{kj}$
the source transition in the opposite direction.  In the neighborhood
of a retained local formula the readout equals
\[
                 T_{k,ji}\circ b_k\circ T_{ij}.
\]
On a slightly smaller compact neighborhood all compositions are
defined.  The source transitions converge smoothly to $T_{ji}$,
so the chain rule gives convergence to
$T_{ji}T_{ij}=\mathrm{id}$ in every derivative.  A finite cover
again makes this uniform on any specified compact subset.

Now adjoin a compact piece in a new chart $q:V_i\to M_\infty$.
Choose compact sets $A,A'$ in the old region and $C\subset V_i$
so that
\[
 P\subset\operatorname{int}A,\qquad
 A\subset\operatorname{int}A',\qquad A'\subset O,
\]
and the new piece lies in $\operatorname{int}C$.
The set $K=C\cap q^{-1}(A')$ is compact and is contained in the
open overlap $q^{-1}(O)$.  This set includes every identification
between the larger old region and the entire closed new piece.
Choose a smooth cutoff $\chi$ equal to one near $K$ with compact
support in $q^{-1}(O)$, and use the actual readout
$r_k=e_{ki}^{-1}F_kq$ there.  Extend
\[
                    a_k(x)=x+\chi(x)(r_k(x)-x)
\]
by the identity off that support.  Smooth convergence of $r_k$ to
the identity gives $\|D(a_k-\mathrm{id})\|_\infty<1/2$ eventually.
The mean value inequality on Euclidean space then gives
$|a_k(x)-a_k(y)|\geq|x-y|/2$.  Thus $a_k$ is injective and locally
invertible; it is proper because it is the identity outside a
compact set, so its open and closed image is all Euclidean space.
It is a smooth diffeomorphism.  For large $k$ its displacement is
smaller than the distance of the support to the chart boundary, so it
preserves $V_i$.  The corrected chart $g_k=e_{ki}\circ a_k$ satisfies
$g_k(y)=F_k(q(y))$ for every $y\in C$ with $q(y)\in A'$.

The margin between $A$ and $A'$ now excludes unwanted identifications.
On the compact product $A\times C$ the distances
$d_{g_k(0)}(F_k(x),g_k(y))$ converge uniformly to
$d_\infty(x,q(y))$.  This follows from the local formulas, the
limits $D_{ij}$, and the small displacement of $a_k$.
All zeros of the limiting distance lie in the open subset
$A\times q^{-1}(\operatorname{int}A')$ of this product:
if $x=q(y)$ and $x\in A$, then $q(y)\in\operatorname{int}A'$.
On its compact complement the limiting distance has a positive
minimum, unless that complement is empty.  Thus any source collision,
for all sufficiently large $k$, has $q(y)\in A'$.  There the
established agreement gives
\[
 F_k(x)=g_k(y)=F_k(q(y)).
\]
Both $x$ and $q(y)$ lie in $O$, where $F_k$ is injective, so
$x=q(y)$.  Conversely that equality implies agreement.  We have
therefore proved the exact equivalence
\[
       F_k(x)=g_k(y)\quad\Longleftrightarrow\quad x=q(y),
       \qquad x\in A,\ y\in C.
\]

Consequently $F_k$ on $\operatorname{int}A$ and $g_kq^{-1}$ on
$q(\operatorname{int}C)$ glue to one injective local diffeomorphism
on their union.  Its local inverses show that it is an open embedding.
This union has compact closure contained in $A\cup q(C)$ and contains
both $P$ and the new compact piece.  On the old part the map is
unchanged; on the new part it has the local formula with correction
$a_k$.  Thus all induction properties, including the exact basepoint
formula, persist.  The old open domain may shrink, but the retained
compact set, enlarged by the new piece, remains inside the next domain.
This distinction is what permits finitely many corrections.

Every compact quotient set has a finite cover by interiors of compact
chart pieces.  The induction therefore gives the required embeddings
on a neighborhood of that set.  Exhaust the connected limit by
relatively compact connected open sets and, for each stage, choose a
source index larger than the previous one for which its embedding is
defined and the first finitely many prescribed jet errors are small.
This diagonal choice gives
Definition~\ref{def:rf-pointed-convergence}.  Embeddings chosen for
different source indices need not agree with one another.  Their
coordinate readouts converge to the identity, so their pullback
metrics have the same limits $G_i$.  This develops the finite patching
step summarized using partitions of unity in
\cite[Theorem 5.6, pp.~86--87]{MT2007}.

\subsection{Boundary escape and completeness at zero}

The chart extraction retains the following coverage property: for each
$R>0$ there are finitely many chart indices and compact coordinate
pieces $C_i\subset V_i$ such that, for all sufficiently large $k$,
\[
                    B_{g_k(0)}(p_k,R)\subset\bigcup_i e_{ki}(C_i).
\]
The pieces and their number are fixed after the subsequence is selected;
pointwise existence of a finite cover on each individual manifold would
not give this assertion.  We now show precisely how this property
enters completeness.

Let $p_\infty=q_{i_0}(p)$ and define
\[
                    \rho(q_i(x))=D_{i_0i}(p,x).
\]
Triangle inequalities show that this is independent of the chosen
representative, continuous and nonnegative, with
$\rho(p_\infty)=0$.  We use it as an auxiliary radius; no identification
with Riemannian distance is needed here.
If $\rho(q_i(x))<R$, then $e_{ki}(x)$ eventually lies in the
covered source $R$-ball.  Some one covering chart occurs along an
infinite subsequence.  In that chart write
$e_{ki}(x)=e_{kj}(y_k)$ with $y_k\in C_j$, and pass to a further
subsequence with $y_k\to y\in C_j$.  Its upper Lipschitz bound gives
\[
 d_{g_k(0)}(e_{ki}(x),e_{kj}(y))
       \leq L_j|y_k-y|\longrightarrow0.
\]
Thus $D_{ij}(x,y)=0$ and $q_i(x)=q_j(y)$.  Every point with
$\rho<R$ belongs to the compact set $\bigcup_jq_j(C_j)$.
For $r<R$ the closed set $\{\rho\leq r\}$ is therefore compact.

The quotient is connected as well.  Otherwise choose a nontrivial
subset that is both open and closed, containing $p_\infty$ and
excluding a point $q$.  Choose $R>\rho(q)$ and split each of the
finitely many covering pieces according to these two quotient subsets.
Both parts are compact.  Any pair from opposite parts has positive
limiting cross-distance, so compactness and uniform convergence make
their source images disjoint for large $k$.  The source basepoint
belongs to the first union and a source representative of $q$ to the
second: a contrary infinite sequence of representations in a wrong
part would give a zero-distance partner there by the preceding
compactness argument.  These two disjoint closed source unions cover
the source $R$-ball and both meet it, contradicting its connectedness.
A Riemannian ball is connected even without completeness: a point
in the ball can be joined to its center by a path of length less than
the radius, and every initial subpath stays in the ball.

Put $E_j$ equal to the component containing $p_\infty$ of
$\{\rho<j+1\}$.  Local connectedness of the manifold makes $E_j$
open.  Its closure is compact and lies in $E_{j+1}$, since it is
connected and lies in $\{\rho<j+2\}$.  The sets $E_j$ exhaust the
quotient: a path from the basepoint to any fixed point has compact
image and hence bounded $\rho$.  Also
\[
                    \rho=j+1\quad\hbox{on }\partial E_j.
\]
Indeed, the boundary of a component of an open set in a locally
connected space is contained in the boundary of that open set;
continuity of $\rho$ then gives the displayed equality.
Use this particular exhaustion in the diagonal embedding construction.

On the compact boundary $\partial E_j$ the source distances
$d_{g_{\sigma(k)}(0)}(p_{\sigma(k)},\phi_k(x))$ converge uniformly
to $\rho(x)$.  To see uniformity, cover that boundary by finitely
many compact chart representatives and combine the original
cross-distance convergence with the local chart approximation of
$\phi_k$.  Given $A>0$, choose $j>A$.  Error less than one now
puts the whole embedded boundary outside the source $A$-ball,
with one eventual index valid for every boundary point.

For large $k$, $\phi_k$ is defined beyond $\overline{E_j}$.
The set $\phi_k(E_j)$ is open and its closure lies in the compact,
hence closed, set $\phi_k(\overline{E_j})$.  A closure point outside
$\phi_k(E_j)$ must be the image of $\partial E_j$, so cannot lie
in the source $A$-ball.  Therefore $\phi_k(E_j)$ is both relatively
open and relatively closed in that connected ball and contains its
center.  It covers the whole ball.

Apply this coverage with radius $2A$.  For any fixed
$x\in B_{g_\infty(0)}(p_\infty,A)$, choose a path of length less
than $A$ from $p_\infty$ to $x$.  Its compact image lies in the
embedding domain eventually, and convergence of lengths puts
$\phi_k(x)$ in the source $2A$-ball.  Coverage supplies
$y\in E_j$ with $\phi_k(y)=\phi_k(x)$.  Both points lie in the
same embedding domain, so injectivity gives $x=y\in E_j$.
Consequently every limit $A$-ball has compact closure.  A Cauchy
sequence is bounded, has a convergent subsequence in one such compact
set, and hence converges.  This proves completeness of $g_\infty(0)$
without assuming completeness of any approximating metric.
Compare \cite[Theorems 5.9--5.11, pp.~88--90; Proposition 5.14,
pp.~90--91; Theorem 5.15, pp.~91--92]{MT2007}.

\subsection{The second time variable}

Curvature and norms converge through the metric jets.  To transfer
condition~2 to a limit ball, choose a short path from $p_\infty$ to
its point at time $s$ and put that path in a compact exhaustion
stage.  Length convergence places its image in a slightly larger
source $s$-ball.  Apply the same source bound at any other time $t$
and then pass to the limit.  Thus, for every $A>0$, some $K_A'<\infty$
bounds limit curvature at all $t\in J$ on every ball
$B_{g_\infty(s)}(p_\infty,A)$, uniformly over $s\in J$.

Here is a direct completeness argument using exactly that quantifier
order.  Fix $t\in J$ and a $g_\infty(t)$-Cauchy sequence $(x_i)$.
It lies in a ball of radius $r$ about $p_\infty$, after increasing
$r>0$.  Condition~2 on the limit gives a curvature bound $K$ on the
fixed set $B_{g_\infty(t)}(p_\infty,3r)$ at every intermediate time
between $t$ and zero.  Hence $|\operatorname{Ric}|\leq n^3K$ there
(this coarse contraction constant suffices).  Pointwise integration
of the metric equation compares tangent norms by the factor
$L=\exp(n^3K|t|)$.

For sufficiently close $x_i,x_j$, a nearly shortest $g_\infty(t)$-path
between them has length less than $r$ and remains in the ball of
radius $3r$.  Comparing its lengths at zero and at $t$ and taking
the infimum gives
$d_{g_\infty(0)}(x_i,x_j)\leq Ld_{g_\infty(t)}(x_i,x_j)$.
The sequence is therefore Cauchy for the complete zero-time metric.
It converges there, hence in the manifold topology, and thus also
for $g_\infty(t)$, which induces the same topology.  This proves
completeness at every interior time.

A curvature bound only on zero-time balls does not supply the fixed
$t$-ball bound in this proof.  The two-time condition is the point of
\cite[Theorem 1.7, p.~179]{Topping2014}; the failure under weaker
hypotheses is discussed in \cite[Theorem 1.3 and Corollary 1.4,
pp.~176--177]{Topping2014}.

\subsection{Ball volumes and static noncollapse in the retained limit}

\begin{lemma}[Volume and noncollapse under complete pointed convergence]
\label{lem:rf-limit-static-noncollapse}
\lean{PoincareMT.PointedGeometricConvergence.eventually_ball_volume_bounds_at,PoincareMT.PointedGeometricConvergence.ball_volume_lower_bound_of_eventually_static_noncollapse}
\uses{def:rf-pointed-convergence,thm:rf-compactness}
Suppose the pointed convergence of
Definition~\ref{def:rf-pointed-convergence} is given in dimension
$n\geq1$.  Fix an interior time $t$ at which the limit metric
$g_\infty=g_\infty(t)$ is complete.  For $x\in M_\infty$,
write $x_i=\phi_i(x)$, $h_i=g_{\sigma(i)}(t)$ and
$\mu_i=\mu_{h_i}$.  Then, for every $r>0$,
\[
       \mu_i(B_{h_i}(x_i,r))
             \longrightarrow\mu_{g_\infty}(B_{g_\infty}(x,r)).
\]
Suppose further that the source metrics are metrically
$\kappa$-noncollapsed, with $\kappa>0$: for every source center
$y$ and radius $r>0$,
\[
 |\operatorname{Rm}_{h_i}|\leq r^{-2}\ \hbox{on }B_{h_i}(y,r)
       \quad\Longrightarrow\quad
       \mu_i(B_{h_i}(y,r))\geq\kappa r^n .
\]
For the conclusion at a given $x$, it suffices that this implication
hold for all sufficiently large $i$ at each fixed radius and
the centers $x_i$.
Then the complete limit metric is metrically
$\kappa$-noncollapsed with the same constant.
\end{lemma}
\begin{proof}
\emph{Coverage at an arbitrary center.}
Fix $C>1$ and $r>0$, and let $R=Cr+1$.
The closed limit $R$-ball about $x$ is compact and lies in one
exhaustion stage.  Uniform convergence of the positive metric
matrices gives, eventually on this ball,
\[
 C^{-1}|v|_{g_\infty}
       \leq |d\phi_i(v)|_{h_i}
       \leq C|v|_{g_\infty}.
\]
The corresponding inverse differential bound is at most $C$.
Take a source path from $x_i$ of length less than $r$.
If it first exited the image of the open limit $R$-ball, the
image of the closed ball, being compact, would contain the
exit point.  The inverse path up to that point has length
less than $Cr<R$, a contradiction.  The inverse endpoint
is in fact at limit distance less than $Cr$.  Hence
\[
 B_{h_i}(x_i,r)\subset
       \phi_i(B_{g_\infty}(x,Cr)).
\]
Conversely, a minimizing limit segment of length less than
$r/C$ lies in the controlled ball.  Its image has length less
than $r$, giving
\[
 \phi_i(B_{g_\infty}(x,r/C))\subset B_{h_i}(x_i,r).
\]
These are statements about the actual embeddings; completeness
of a source metric is unnecessary.

\emph{The volume squeeze at every radius.}
The singular values of $d\phi_i$ with respect to these metrics
lie between $C^{-1}$ and $C$, so its volume Jacobian lies between
$C^{-n}$ and $C^n$.  Change of variables on the embedding image
and the two inclusions give
\[
 C^{-n}\mu_{g_\infty}(B(x,r/C))
 \leq\mu_i(B(x_i,r))
 \leq C^n\mu_{g_\infty}(B(x,Cr)).
\]
To let $C\downarrow1$, we need continuity in radius, including
radii meeting the cut locus.  A positive-radius intrinsic sphere
in a complete manifold is covered by the exponential image of
the Euclidean sphere of the same radius: choose a minimizing
segment to each of its points.  That Euclidean sphere has
$n$-dimensional Lebesgue measure zero.  The exponential map is
smooth; on local coordinate neighborhoods it is Lipschitz on
compact subsets and therefore preserves such null sets.
Riemannian volume has a smooth positive coordinate density.
A finite cover of the compact tangent sphere thus shows that
the intrinsic sphere has zero volume.

For radii tending to $r$, the indicators of the corresponding
balls converge off this null sphere.  They are eventually
dominated by the indicator of the $(r+1)$-ball, whose volume is
finite by completeness and local finiteness.  Dominated
convergence proves continuity of ball volume at $r$.
The displayed squeeze now proves the asserted volume limit.

\emph{Strict curvature margins.}
Assume $|\operatorname{Rm}_{g_\infty}|\leq r^{-2}$ on $B(x,r)$.
Choose $0<\rho<a<r$ and put $C=a/\rho>1$.
The compact closure of $B(x,a)$ lies inside $B(x,r)$.
Metric convergence through order two, uniform positive
definiteness, and the coordinate formulas
$\Gamma=g^{-1}*\partial g$ and
$\operatorname{Rm}=g*(\partial\Gamma+\Gamma*\Gamma)$
give uniform curvature-norm convergence on that closure.
Here the star denotes a finite sum of tensor contractions;
only inverses of uniformly positive matrices and their first
two derivatives occur.  A finite coordinate cover makes
the convergence uniform even when the point varies.
Because $r^{-2}<\rho^{-2}$, all sufficiently late embedded
points of that compact set have curvature norm less than
$\rho^{-2}$.

Coverage at radius $\rho$ places the entire source $\rho$-ball
inside the image of $B(x,a)$.  Thus the source noncollapse
implication applies and gives
$\mu_i(B(x_i,\rho))\geq\kappa\rho^n$.
Volume convergence at $\rho$ yields
$\mu_{g_\infty}(B(x,\rho))\geq\kappa\rho^n$.
This ball is contained in $B(x,r)$; let $\rho\uparrow r$ to
obtain the required $\kappa r^n$ lower bound.
The eventual index may depend on $x,\rho,a$; no common index
over all radii or centers was used.
\end{proof}

Applying this lemma separately at each complete interior slice
preserves static noncollapse throughout a pointed flow limit.
A curvature bound on a backward ball cylinder includes its
terminal spatial bound, so static noncollapse at every slice
also gives parabolic noncollapse with the same constant.
This is the complete-ball passage needed in the compactness
arguments of \cite[Corollary 44.1, pp. 2682--2683, and
Appendix E, pp. 2848--2849]{KleinerLott2013}.

\section{Inputs to the later arguments}
\label{sec:rf-interfaces}

The later surgery construction starts from
Proposition~\ref{prop:rf-normalize} and
Theorem~\ref{thm:rf-local}.  On each smooth interval it uses the
scalar clock $-6/(1+4t)$ and the pinching condition with the same
absolute time; a restart at time $a$ does not replace $t$ by $t-a$
inside the logarithmic barrier.  Scalar upper bounds then give full
curvature bounds by \eqref{eq:rf-full-bound}.

Reduced geometry and evolving-curve estimates use
Proposition~\ref{prop:rf-evolution},
Theorem~\ref{thm:rf-shi}, and
Proposition~\ref{prop:rf-metric-comparison} on the same component,
metric, and time slab as their application.  Ancient-model limits
use Theorem~\ref{thm:rf-compactness}; they must verify its volume
normalization, noncollapse and two-time bounds.  They may use the
bounded ancient Harnack consequence before the asymptotic
volume-ratio argument.  Use of the full slice-bounded
Theorem~\ref{thm:rf-harnack} requires the additional volume-to-curvature
argument in Section~\ref{sec:rf-harnack}.
